% concatenated from main.tex
\documentclass{article}


\usepackage{amssymb}
\usepackage{amsmath}
\usepackage{amsthm, amsfonts, mathrsfs}
\usepackage{xcolor}

\usepackage{enumitem}

\usepackage{float}

% This package allows us to use symbols from `linear logic' - in our context, the symbols for operations on relations 
\usepackage{cmll}

\usepackage{tikz}

\usetikzlibrary{calc}




\newtheorem{thm}{Theorem}[section]
\newtheorem{cor}[thm]{Corollary}
\newtheorem{claim}{Claim}
\newtheorem{prop}[thm]{Proposition}
\newtheorem{prop1}[thm]{Property}
\newtheorem{lem}[thm]{Lemma}
\newtheorem{re}[thm]{Remark}

\newtheorem{ex}[thm]{Example}
\newtheorem{df}{Definition}
\newtheorem{conj}{Conjecture}
\newtheorem{Fact}[thm]{Fact}


\newtheorem{qu}{Question}
\newtheorem{pr}[qu]{Problem}
\renewcommand*{\thequ}{\Alph{qu}}

\newcommand{\comment}[1]{}

\newcommand{\lb}{\left\{}
\newcommand{\rb}{\right\}}
\newcommand{\la}{\left<}
\newcommand{\ra}{\right>}

\DeclareMathOperator{\fini}{fi}

\newcommand{\tq}{\ge_T}
\newcommand{\te}{=_T}

\newcommand{\gtq}{\ge_{T}}
\newcommand{\gte}{=_{T}}

\newcommand{\K}{\mathcal{K}}
\newcommand{\F}{\mathcal{F}}

\newcommand{\omom}{\omega^\omega}

\newcommand{\cl}[1]{\overline{#1}}

\DeclareMathOperator{\ci}{c}
\DeclareMathOperator{\nci}{\neg c}

\DeclareMathOperator{\ii}{ii}

\DeclareMathOperator{\cof}{cof}
\DeclareMathOperator{\add}{add}




\title{The Cofinality of Generating Families}
\author{Paul Gartside and Thomas Gilton}
%\date{September 2025}

\begin{document}

\maketitle
\begin{abstract}
The topology of a separable metrizable space $M$ is \emph{generated} by a family $\mathcal{C}$ of its subsets provided that a set $A\subseteq M$ is closed in $M$ if and only if $A\cap C$ is closed in $C$ for each $C\in \mathcal{C}$. The \emph{sequentiality number}, $\mathop{seq}(M)$, and \emph{$k$-ness number}, $\mathop{k}(M)$,   of $M$, are the minimum size of a generating family of convergent sequences, respectively compact subsets. 

Let $\mathfrak{b}$ be the minimum size of an unbounded set in $\omega^\omega$ with the mod finite order. For a cardinal $\kappa$, the \emph{covering number}, $\mathop{cov}(\kappa)$, is the minimum size of a family of countable subsets of $\kappa$ so that every countable subset of $\kappa$ is contained in an element of the family. It is shown using the Tukey order on relations that (1) $\mathop{seq}(M)=\mathop{cov}(|M|)\cdot \mathfrak{b}$,  unless $M$ is locally small (every point of $M$ has a neighborhood of size strictly less than $|M|$) in which case $\mathop{seq}(M)=\lim_{\mu <|M|} \mathop{cov}(\mu)\cdot \mathfrak{b}$ and 
(2) $k(M)$ is in the interval $[kc(M)\cdot\mathfrak{b},\mathop{cov}(kc(M))\cdot \mathfrak{b}]$, where $kc(M)$ is the minimum number of compact sets that cover $M$. 

Solutions to problems of van Douwen's on the $k$-ness number of analytic and of co-analytic spaces are deduced. 

\smallskip
Keywords: separable metrizable, sequential, $k$-space, cardinal characteristics of the continuum, small cardinals, cardinal invariants,  Tukey order, \textsc{PCF} theory. 

MSC Classification: 03E04, 03E15, 03E17, 06A07, 54A25, 54D45, 54D50,  54D55, 54E35. 
\end{abstract}



\section{Introduction}
For any subset $M$ of the real line its convergent sequences determine which subsets of $M$ are closed,   hence the convergent sequences determine the topology, and hence the closure of any set and which real valued functions on $M$ are continuous. It is natural to ask \emph{how many sequences suffice to generate the topology?}  
More concretely, the topology of a space $M$ (all spaces in this paper are separable and metrizable)   is \emph{generated} by a family $\mathcal{C}$ of its subsets provided that a subset $A$ of $M$ is closed in $M$ if and only if $A\cap C$ is closed in $C$ for each $C\in \mathcal{C}$. Then $\mathop{seq}(M)$, the \emph{sequentiality number} of $M$, is the minimum size of a family of convergent sequences in $M$ generating the topology of $M$. 
It is also natural to define  $\mathop{k}(M)$, the \emph{$k$-ness number} of $M$, as the minimum size of a generating family  of compact sets. In this paper we compute $\mathop{seq}(M)$ and $k(M)$ in terms of intrinsic properties of the separable metrizable space $M$.

Van Douwen introduced the $k$-ness number  
in his survey article, \cite{vanDou}, on cardinal characteristics of the continuum and small cardinals arising in topology. Observe that, for any separable metrizable space  $M$, we have $k(M) \le \mathop{seq}(M) \le \mathfrak{c}$, the continuum ($\mathfrak{c}=|\mathbb{R}|=2^{\aleph_0}$). Van Douwen noted that $k(M) \le \omega$ only when $M$ is locally compact, and he showed that when  $M$ is not locally compact we have $\mathfrak{b} \le k(M)$, where $\mathfrak{b}$ is the \emph{bounding number} (minimum size of an unbounded set in $\omega^\omega$ with the mod finite order, $<^*$). 
Van Douwen computed the $k$-ness number for a handful of spaces, notably the rationals, but left open (see \cite[Question~8.11]{vanDou}) the computation of $k(M)$ for $M$ analytic or co-analytic. We conclude the paper with the solution to van Douwen's question. 

Our results depend on two sources of techniques: first the theory of relations, specifically the Tukey order,   and second Shelah's \textsc{PCF}  theory. 
We now give an overview of how these techniques combine to enable us to compute the sequentiality and $k$-ness numbers. It is convenient to do so initially for $\mathop{seq}(M)$, where our results are essentially complete (up to the limits of our understanding of \textsc{PCF} theory); and then for $k(M)$, where our results loosely mirror those for sequentiality.

From work by the first author and Ziqin Feng in \cite{FG_Shape} we have a detailed understanding, mediated by the Tukey order, of the relations $\mathbf{seq}(M)$ and $\mathbf{k}(M)$,  whose cofinal sets are generating families of convergent sequences and compact subsets, respectively, and hence whose cofinalities (minimum size of a cofinal set) are $\mathop{seq}(M)$ and $k(M)$, respectively. 
 A map $\phi_+$ from one relation, $\mathbf{A}$, to another, $\mathbf{B}$, is a Tukey morphism if it carries cofinal sets to cofinal sets. The Tukey order is given by $\mathbf{A} \tq \mathbf{B}$ if there is a Tukey morphism of the first relation to the second. If $\mathbf{A} \tq \mathbf{B}$ then the cofinality of $\mathbf{A}$ is at least as big as the cofinality of $\mathbf{B}$. Consequently the Tukey order is a convenient tool to compare and compute cofinalities of relations such as $\mathbf{seq}(M)$ and $\mathbf{k}(M)$. 
(See Section~\ref{ssec:prelim_rel} for background material on relations and Tukey morphisms.) 

The outcome for the sequentiality number of a separable metrizable space $M$ is that $\mathop{seq}(M)=\mathop{sam}(|M|)\cdot \mathfrak{b}$,  unless $M$ is locally small (every point of $M$ has a neighborhood of size strictly less than $|M|$) in which case $\mathop{seq}(M)=\lim_{\mu <|M|} \mathop{sam}(\mu)\cdot \mathfrak{b}$. 
Here, for any cardinal $\kappa$, the \emph{sampling number} of $\kappa$ is the least number of countable subsets of $\kappa$ needed to have infinite intersection with every countably-infinite subset of $\kappa$.  Evidently the sampling number is a purely set-theoretic object. It is easily seen that $\mathop{seq}$ is ranked by the size of the space: if $|M| \le |N|$ then $\mathop{seq}(M) \le \mathop{seq}(N)$. Also among all spaces $M$ of a fixed cardinality, $\mathop{seq}(M)$ can only take on at most two values, and can attain more than one only when $|M|$ has countable cofinality. 

To compute the sequentiality number, then, we need to understand the sampling number and its relation to the bounding number, $\mathfrak{b}$, and the continuum, $\mathfrak{c}$ (because separable metrizable spaces have cardinality no more than $\mathfrak{c}$).  
One of the fundamental objects of \textsc{PCF} theory is the \emph{covering number} of a cardinal $\kappa$ which is the minimal size of a family  of countable subsets of $\kappa$ such that every countable subset is contained in a member of the family. 
Clearly the sampling and covering numbers are closely related. In fact, see Theorem~\ref{th:cov=sam}, despite the formal difference between them, they coincide: for every $\kappa$ we have $\mathop{sam}(\kappa)=\mathop{cov}(\kappa)$. 
From \textsc{PCF} theory much is known about the covering number and we recapitulate this next. However 
our results above on the sequentiality number, as well as those below on the $k$-ness number, are most naturally established for the sampling number. 
Since the sampling and covering numbers are equal, we can substitute one for the other at our convenience. 
For consistency we have stated our results below using the covering number, but in proofs often use the sampling number instead without explicitly (and repeatedly) stating the equality.  The concluding paragraph of Section~\ref{ssec:seqsam} gives details on how our understanding of the covering number informs our knowledge of the sequentiality number.

It is known that $\kappa \le \mathop{cov}(\kappa)$, but $\mathop{cov}(\kappa) \ge \kappa^+$ for uncountable $\kappa$ with countable cofinality. Further, it is a basic tenet of \textsc{PCF} theory that the covering number is `tame' - $\mathop{cov}(\kappa)$ does not jump up much above $\kappa^+$. 
Indeed, at least assuming the Continuum Hypothesis, \textsc{CH},  large cardinal strength is required  to get $\mathop{cov}(\kappa) > \kappa^+$. However, it is also known that, assuming the existence of a supercompact cardinal, it is consistent that \textsc{CH} holds and  $\mathop{cov}(\kappa) > \kappa^+$ for $\kappa=\aleph_\omega$. 
Applying random real forcing to the latter model,  allows us to construct, again modulo large cardinals, models where, for example,  $\mathfrak{b}=\aleph_1$ (and so is a non-factor in the possible values of $\mathop{seq}$), $\mathop{cov}(\aleph_\omega)>\aleph_{\omega+1}$ and $\mathfrak{c}$ is as large as we like. Section~\ref{ssec:pcf} provides  details on the above, and concludes with a discussion of the cofinality of products of countable sets of cardinals with the mod finite order, and in particular the existence of mod finite scales. 
This enables the construction of models, modulo large cardinals, in which we can control the covering numbers of families of cardinals. 


Finally, in Section~\ref{ssec:ksam} we turn to the question of how many compact subsets suffice to generate the topology of a separable metrizable space. 
Let $kc(M)$, the \emph{compact covering number}, be the minimum number of compact subsets needed to cover a separable metrizable $M$. 
The case of $\sigma$-compact spaces (i.e. $kc(M) \le \omega$) was known to van Douwen, so we move on to  $M$ with $kc(M)$ uncountable.  
For such $M$ it is no surprise that $kc(M) \le k(M)$, but it was a surprise to the authors that $k(M) \le \mathop{cov}(kc(M))\cdot \mathfrak{b}$, and so - as discussed above in  the context of the sequentiality number - the compact covering and $k$-ness numbers are tightly coupled. But less so than the sequentiality number of a space and its cardinality. 
Recall for a fixed $\kappa$, $\mathop{seq}(M)$, for  any $M$ with $|M|=\kappa$, is constrained to at most two values. However we show, modulo large cardinals, for every finite $K\ge 2$  it is consistent that,  for a fixed $\kappa$, $k(M)$, for  any $M$ with $kc(M)=\kappa$, takes on exactly $K$ values, and it is consistent that,  for a fixed $\kappa$, there are infinitely many values of $k(M)$, for   $M$ with $kc(M)=\kappa$. 
Section~\ref{ssec:ksam} concludes with the solution to van Douwen's questions on the $k$-ness numbers of analytic and co-analytic spaces.

%\paragraph{More on PCF and Open Questions} 
%The connection between \textsc{PCF} theory and the sampling, sequentiality and $k$-ness numbers seems to the authors to be novel and raises interesting questions. 


%The sequentiality and $k$-ness numbers depend on the coveringing number of a cardinal $\kappa$ where $\kappa \le \mathfrak{c}$. Historically, however, \textsc{PCF} theory has focused - motivated by the Singular Cardinals Hypothesis, for  example - on singular cardinals in the presence of \textsc{CH}. So there is a mismatch between our needs and the literature. In particular, as mentioned above,  assuming  \textsc{CH}, it is known that large cardinal strength is required  to get $\mathop{cov}(\kappa) > \kappa^+$. 
%Can the Continuum Hypothesis be removed? See Question~\ref{qu:large_card}. 

%Third, the lower bounds on the sampling number come from the additivity and existence of scales on products of cardinals with the mod finite order. \textsc{PCF} theory gets its very name from the cofinalities of products ordered modulo an ultrafilter, about which a great deal is known. Again there is a mismatch - mod finite \textsl{versus} modulo an ultrafilter. See Question~\ref{qu:mod_finite} for details, and see also \cite{Gilton_PCF_Tukey} for related problems. 
%Because the covering number of the first fixed point of the $\aleph$-function can be arbitrarily large, lower bounds of the sampling number are of special interest in this case, see Question~\ref{qu:samUnBounded}.


\section{Preliminaries}
\subsection{Relations and Tukey Order}\label{ssec:prelim_rel} 
\paragraph{Relations, Cofinality, Morphisms}
A \emph{relation} is a triple $\mathbf{A}=(A_-,A_+,A)$, where $A_-$ is the domain, $A_+$ is the co-domain, and $A$ is a collection of pairs  (a \emph{general relation}). 
Write $xAy$ if $(x,y) \in A$. 
Abbreviate to $(A',A)$  a relation $(A_-,A_+, A)$  where $A_-=A'=A_+$. 
A subset $C$ of $A_+$ is \emph{cofinal} in $\mathbf{A}$ if for every $x$ from $A_-$ there is a $y$ from $C$ such that $xAy$, and  the \emph{cofinality} of $\mathbf{A}$, denoted $\cof(\mathbf{A})$, is the minimum size of a cofinal set. 
A subset $S$ of $A_-$ is \emph{bounded} if for some $y$ in $A_+$, for all $x$ from $S$ we have $xAy$. 
The \emph{additivity} of a relation $\mathbf{A}$, denoted $\mathop{add}(\mathbf{A})$, is the minimum size of an unbounded set.  
Equivalently, it is the cofinality of the dual $\mathbf{A}^\perp=(A_+,A_-,\neg A^{-1})$ of $\mathbf{A}$.

A \emph{Tukey morphism} from $\mathbf{A}$ to $\mathbf{B}=(B_-,B_+,B)$ 
was  originally  defined (see \cite{Vojtas}), to be a pair $( \psi_-,\phi_+)$ of functions such that (i) $\phi_+: A_+ \to B_+$, (ii) $ \psi_-: B_- \to A_-$ and (iii) $\psi_-(b) A a \implies b B \phi_+(a)$, where  $b \in B_- \ \text{and}\ a \in A_+$. 
Note that the map $\phi_+$ of a Tukey morphism pair 
carries cofinal sets to cofinal sets.  
It can be shown conversely that if $\phi_+:A_+ \to B_+$  carries cofinal sets to cofinal sets then there exists a $\psi_-$ so that $(\psi_-,\phi_+)$ is a Tukey morphism. Consequently we also call a function $\phi_+:A_+ \to B_+$  mapping cofinal sets to cofinal sets a Tukey morphism. 
If such a morphism exists we write $\mathbf{A} \gtq \mathbf{B}$.  
The relation $\gtq$, the \emph{Tukey order} on relations,  is transitive, hence $\gte$, defined by $\mathbf{A} \gte \mathbf{B}$ if and only if $\mathbf{A} \gtq \mathbf{B}$ and $\mathbf{B} \gtq \mathbf{A}$, is an equivalence relation,  \emph{Tukey equivalence} on relations, whose equivalence classes are called \emph{Tukey types}.

Suppose $(\psi_-,\phi_+)$ is a Tukey morphism pair from $\mathbf{A}$ to $\mathbf{B}$. 
Then $\phi_+$ maps cofinal sets to cofinal sets, and so $\cof(\mathbf{A}) \ge \cof(\mathbf{B})$. 
One can check that $(\phi_+,\psi_-)$ is a Tukey morphism pair from $\mathbf{B}^\perp$ to $\mathbf{A}^\perp$, and so $\mathop{add}(\mathbf{A}) \le \mathop{add}(\mathbf{B})$.

For any set $S$, equality on $S$ gives a relation, $(S,S,=)=(S,=)$. Note $(S,=) \te (S',=)$ if and only if $|S|=|S'|$. 
If $\mathcal{S}$ is a family of sets then $(S,\mathcal{S},\in)$ is a relation whose cofinal sets are the covers of $S$ from $\mathcal{S}$. 
On $\omega^\omega$ let $<^*$ be the mod finite order, so $f<^*g$ if $f(n)<g(n)$ except for finitely many $n$. Then the \emph{dominating number}, $\mathfrak{d}$, is $\cof(\omega^\omega,<^*)$; while the \emph{bounding number}, $\mathfrak{b}$, is the additivity of $(\omega^\omega,<^*)$. 

\paragraph{Operations on Relations}  
Given relations $\mathbf{A}$ and $\mathbf{B}$ define 
$\mathbf{A} \times \mathbf{B}$ to be the standard product, $(A_-\times B_-,A_+\times B_+,A\times B)$, and 
$\mathbf{A} + \mathbf{B}$,  to be $(A_- \oplus B_-,A_+ \oplus B_+,S)$ where $x S y$ if and only if either $x \in A_-$, $y \in A_+$ and $xAy$ or $x\in B_-$, $y\in B_+$ and $xBy$. 
More generally, if $\{\mathbf{A}_\lambda=((A_\lambda)_-,(A_\lambda)_+,A_\lambda) : \lambda \in \Lambda\}$ is a family of relations then its sum is $\sum_{\lambda \in \Lambda} \mathbf{A}_\lambda = (\bigoplus_{\lambda \in \Lambda} (A_\lambda)_-,\bigoplus_{\lambda \in \Lambda} (A_\lambda)_+, S)$ where $x S y$ if and only if for some $\lambda$, $x\in (A_\lambda)_-, y \in (A_\lambda)_+$ and $xA_\lambda y$.
We need the following facts about the sum operation:
\begin{enumerate}[noitemsep,topsep=1pt]
    \item[(A)] If for all $\lambda\in\Lambda$, $\mathbf{A}_\lambda = \mathbf{A}$, then 
    $
    \sum_{\lambda \in \Lambda} \mathbf{A}_\lambda \gte \mathbf{A} \times (|\Lambda|,=)$. 
    \item[(B)]  If  $\Lambda'$ is a subset of $\Lambda$ then 
    $
    \sum_{\lambda \in \Lambda} \mathbf{A}_\lambda \gtq \sum_{\lambda \in \Lambda'} \mathbf{A}_\lambda$. 
    \item[(C)]  If for every $\lambda$ in $\Lambda$,  $\mathbf{A}_\lambda \gtq \mathbf{B}_\lambda$, then 
    $
    \sum_{\lambda \in \Lambda} \mathbf{A}_\lambda \gtq \sum_{\lambda \in \Lambda} \mathbf{B}_\lambda$. 
\end{enumerate}

\paragraph{Some Special Relations} 
Let $(S_n)_n$ be a sequence of sets. Let $\left(\prod_n [S_n]^{<\omega}\right)_\infty$ be all $(F_n)_n$ in the product such that $F_n\ne \emptyset$ for infinitely many $n$.
Define a general relation $i_\infty$ on sequences of sets by $(F_n)_n i_\infty (G_n)_n$ if and only if for infinitely many $n$ we have $F_n \cap G_n \ne \emptyset$. 
When all the $S_n$ are equal, say to $S$, then we write $\left([S]^{<\omega}\right)^{\omega}_\infty$ for $\left(\prod_n [S_n]^{<\omega}\right)_\infty$. 
For a cardinal $\kappa$, we now define the three relations $P(\kappa)$, $Q(\kappa)$, and $S(\kappa)$. Let $P(\kappa)=(\left([\kappa]^{<\omega}\right)^{\omega}_\infty,i_\infty)$. 
If $\cof(\kappa) \ne \omega$ then set $Q(\kappa)$ and $S(\kappa)$ to be the trivial relation. Otherwise, when $\cof(\kappa)=\omega$, fix a strictly  increasing sequence of infinite cardinals, $(\kappa_n)_n$ with limit $\kappa$, and set $Q(\kappa)=(\left(\prod_n [\kappa_n]^{<\omega}\right)_\infty,i_\infty)$ and  $S(\kappa)=\sum_n P(\kappa_n) \times (\kappa_n,=)$. 
From \cite[Lemma~2.14]{FG_Shape} we know $Q(\kappa)$ and $S(\kappa)$ are well-defined in the case when $\kappa$ has countable cofinality, because their  Tukey type is independent of the choice of cofinal sequence, $(\kappa_n)_n$.
We know \cite[Lemma~2.10]{FG_Shape} that $P(\omega) \te (\omom,\le_\infty)$, where $f \le_\infty g$ if $\{n \in \omega : f(n) \le g(n)\}$ is infinite, and hence $\cof(P(\omega))=\mathfrak{b}$. We end this subsection with the following simple lemma
\begin{lem}\label{lem:cofatleastkappa}
    $\cof(P(\kappa))$ is at least $\kappa$. 
\end{lem}
    \begin{proof} Suppose that $\mathcal{C}' \subseteq ([\kappa]^{<\omega})^{\omega}_{\infty}$ has size strictly less than $\kappa$. Then the set $S=\bigcup \{ \bigcup_n F_n : (F_n)_n \in \mathcal{C}'\}$ has size strictly less than $\kappa$, so pick $\alpha \in \kappa \setminus S$. 
Now we see that $(\{\alpha\})_n$ witnesses that $\mathcal{C}'$ is not cofinal.
\end{proof}




\subsection{\textsc{PCF} Theory}\label{ssec:pcf} 

%Now we draw some connections between the current discussion and Shelah's \textsc{PCF} theory. 
Two elements of  Shelah's \textsc{PCF} theory 
play a key role in our results, namely the covering number and mod finite scales.

\paragraph{The Covering Number} 
Relations of the form $([\mu]^{\delta},\subseteq)$, and their cofinalities, where $\mu$ is singular and  $\delta\leq\mu$, are central objects of \textsc{PCF} theory. For the purposes of this paper, the case $\delta=\omega$ is the most important, and, simplifying terminology and notation, we define the \emph{covering number}, $\mathop{cov}(\kappa)$, of a cardinal $\kappa$, to be $\cof([\kappa]^\omega,\subseteq)$. Basic facts about the covering number are recorded below, see \cite{ShelahBook}.

\begin{lem}\label{lem:samIncreasing}
    If $\kappa\geq\lambda$, then $\mathop{cov}(\kappa)\geq\mathop{cov}(\lambda)$.
\end{lem}



\begin{lem} \label{l:sam1} 
Let $\kappa$ be an uncountable cardinal. Then
\begin{enumerate}[noitemsep,topsep=1pt]
\item[(i)]  $\kappa \le \mathop{cov}(\kappa) \le \kappa^\omega$, and $\mathop{cov}(\kappa^\omega)=\kappa^\omega$, 

\item[(ii)] if $\cof(\kappa)=\omega$ then $\mathop{cov}(\kappa) > \kappa$,

\item[(iii)] if $\kappa=\mu^+$ then     $\mathop{cov}(\kappa) = \mathop{cov}(\mu)\cdot\kappa$,  and

\item[(iv)] if $\kappa$ is a limit and $\cof(\kappa)>\omega$  then $\mathop{cov}(\kappa)=\lim \{\mathop{cov}(\mu) : \mu<\kappa\}$. 
\end{enumerate}
\end{lem}

\begin{lem}\label{l:sam_aleph_n}
    For every $n \in \mathbb{N}$ we have $\mathop{cov}(\aleph_n)=\aleph_n$.
\end{lem} 




In several cases, \textsc{PCF} theory places strong upper bounds on the possible values of covering numbers.
The following proposition, which is surely folklore, gives an example.



\begin{prop}\label{pr:fixpt}
    Let $\lambda$ be the least fixed point of the $\aleph$-function. 
    Let $\kappa<\lambda$ be infinite, and write $\kappa=\aleph_\alpha$. 
    Then $\mathop{cov}(\kappa)=\cof([\kappa]^\omega,\subseteq)<\aleph_{|\alpha|^{+4}}$.
\end{prop}
\begin{proof}
    We will need several facts, mostly from \textsc{PCF} theory, which we list here and which can be found in a standard \textsc{PCF} reference, such as \cite{AbrahamMagidor}:
    \begin{enumerate}[nosep,topsep=1pt]
        \item[(1)] 
        $\cof([\mu]^{\nu_1},\subseteq)\leq\cof([\mu]^{\nu_2},\subseteq)\cdot\cof([\nu_2]^{\nu_1},\subseteq)$, when
        $\mu$ is infinite and $\nu_1\leq\nu_2$. 
        \item[(2)] Suppose that $\mu$ is a singular cardinal and $\kappa<\mu$ is any infinite cardinal so that the interval $A:=(\kappa,\mu)\cap\textsc{REG}$ of regular cardinals has size $\leq\kappa$.
        Then 
        $\cof([\mu]^{\kappa},\subseteq)=\operatorname{maxpcf}(A)$.
        \item[(3)] Suppose that $A$ is a progressive interval of regular cardinals.
        Then $\operatorname{pcf}(A)$ is an interval of regular cardinals.
        \item[(4)] If $A$ is a progressive interval of regular cardinals, then 
        $|\operatorname{pcf}(A)|<|A|^{+4}$.        
    \end{enumerate}

We prove the proposition by induction on $\kappa\geq\omega$. The result is trivial for $\kappa=\omega$, since the relevant cofinality is simply equal to 1.
We next take care of the successor case. 
Suppose that $\kappa=\aleph_{\alpha+1}$, for some $\alpha<\lambda$ and that the result is known for $\aleph_\alpha$.
Then
$
\cof([\aleph_{\alpha+1}]^\omega,\subseteq)=\cof([\aleph_\alpha]^\omega,\subseteq)\cdot\aleph_{\alpha+1}
$.
By induction, $\cof([\aleph_\alpha]^\omega,\subseteq)<\aleph_{|\alpha|^{+4}}$.
We also trivially have $\alpha+1<|\alpha|^{+4}$.
Thus both terms on the right-hand side are below $\aleph_{|\alpha|^{+4}}\leq\aleph_{|\alpha+1|^{+4}}$, which completes the proof in the successor case.

We now address the limit case.
Suppose that $\kappa=\aleph_\alpha$ where $\alpha$ is a limit.
To structure the following argument, we begin with an application of fact (1) above, namely that 
$(\ast)\ \cof([\kappa]^\omega,\subseteq)\leq\cof([\kappa]^{|\alpha|},\subseteq)\cdot\cof([|\alpha|]^\omega,\subseteq)
$.
Our argument will proceed by bounding both of the terms on the right-hand side of $(\ast)$.

To begin, note that $\alpha<\aleph_\alpha$, since $\kappa<\lambda$ and $\lambda$ is the least fixed point of the $\aleph$-function.
Since $\alpha$ is a limit, $|\alpha|$ is infinite.
Also, the interval $A:=(|\alpha|,\kappa)\cap\textsc{REG}$ of regular cardinals has size $|\alpha|$: indeed, as $\kappa=\aleph_\alpha$, the function taking each $\nu\in A$ to its index $\beta$ in the $\aleph$-sequence is an injection of $A$ into $\alpha$.
By fact (2) above, we conclude that $\cof([\kappa]^{|\alpha|},\subseteq)=\operatorname{maxpcf}(A)$.
As $A$ is a progressive (since $|A|=|\alpha|<\min(A)$) interval of regular cardinals, we conclude from facts (3) and (4) that $\operatorname{pcf}(A)$ is an interval of regular cardinals of size $<|A|^{+4}=|\alpha|^{+4}$.

We now claim that $\operatorname{maxpcf}(A)<\aleph_{|\alpha|^{+4}}$.
Suppose otherwise, for the sake of a contradiction.
Certainly $\operatorname{pcf}(A)$ contains $\min(A)$ and $\operatorname{maxpcf}(A)$.
Since $\operatorname{pcf}(A)$ is an interval of regular cardinals, it therefore contains the interval $[\min(A),\operatorname{maxpcf}(A)]\cap\text{REG}$.
Since $\aleph_{|\alpha|^{+4}}<\operatorname{maxpcf}(A)$, $\operatorname{pcf}(A)$ therefore contains the interval $[\min(A),\aleph_{|\alpha|^{+4}})\cap\text{REG}$ of regular cardinals.
As $\min(A)=|\alpha|^+<\aleph_{|\alpha|^{+4}}$, we conclude that $\operatorname{pcf}(A)$ has at least $|\alpha|^{+4}$-many elements, contradicting the fact that $|\operatorname{pcf}(A)|<|\alpha|^{+4}$.

Bound the first term on the right-hand side of $(\ast)$, by stringing together the previous results to see 
$\cof([\kappa]^{|\alpha|},\subseteq)=\operatorname{maxpcf}(A)<\aleph_{|\alpha|^{+4}}$.

To bound the second term, we use the induction hypothesis.
Write $|\alpha|=\aleph_\beta$ for some $\beta<|\alpha|$. 
By induction $\cof([|\alpha|]^\omega,\subseteq)<\aleph_{|\beta|^{+4}}$.
Since $\beta<|\alpha|$ (because $|\alpha|$ is not a fixed point of the $\aleph$-function), we know that $|\beta|^{+4}<|\alpha|^{+4}$.
Thus $\cof([|\alpha|]^\omega,\subseteq)<\aleph_{|\alpha|^{+4}}$.

We  conclude that both terms on the right-hand side of $(\ast)$ are $<\aleph_{|\alpha|^{+4}}$, and this completes the induction.
\end{proof}

Here we note that the situation is more complex at fixed points of the $\aleph$-function. Assuming the consistency of certain large cardinal assumptions, $\cof([\kappa]^\omega,\subseteq)$ can be arbitrarily large when $\kappa$ is the first fixed point of the $\aleph$-function; this result follows almost immediately from the following theorem, due to Gitik:

\begin{thm} (Gitik, \cite{Gitik:noBounds}) Let $\kappa$ be a cardinal of cofinality $\omega$ with certain large cardinals below $\kappa$.\footnote{Namely, for each $\tau<\kappa$, the set $\lb\alpha<\kappa:o(\alpha)\geq\alpha^{+\tau}\rb$ is unbounded in $\kappa$; note that $o(\alpha)$ is the Mitchell order of a cardinal $\alpha$.} Let $\lambda>\kappa$ be any cardinal. Then there is a forcing extension in which $\kappa$ is the first fixed point of the $\aleph$-function, $2^\kappa\geq\lambda$, and the \textsc{GCH} holds below $\kappa$.
\end{thm}

\begin{cor} (Gitik, \cite{Gitik:noBounds})\label{cor:noBounds}
    Assuming the existence of certain large cardinals, $\mathop{cov}(\kappa)$ can be arbitrarily large, where $\kappa$ is the first fixed point of the $\aleph$-function, and \textsc{CH} holds.
\end{cor}
\begin{proof}
    Work in Gitik's model from above where $\kappa$ is the first fixed point of the $\aleph$-function, \textsc{GCH} holds below $\kappa$, and $2^\kappa\geq\lambda$ (where $\lambda>\kappa$ is arbitrary). 
    We first note that $\kappa^\omega=\cof([\kappa]^\omega,\subseteq)\cdot2^{\aleph_0}=\cof([\kappa]^\omega,\subseteq)$, using the fact that $2^{\aleph_0}=\aleph_1$ in the model under consideration. Next observe that $\kappa^\omega=2^\kappa$, using the \textsc{GCH} below $\kappa$. Putting these together, we conclude that $\cof([\kappa]^\omega,\subseteq)=2^\kappa\geq\lambda$.
\end{proof}



Optimal upper bounds follow from `no large cardinal' hypotheses. 
Indeed, by work of Gitik \cite{Gitik:Strength}, we know that if there does not exist a cardinal $\kappa$ with Mitchell order $\kappa^{++}$ (written $o(\kappa)=\kappa^{++})$, and if $\mu$ is a singular cardinal of cofinality $\omega$, \emph{and} if $\mathfrak{c}=2^\omega<\mu$ (note the cardinal arithmetic assumption here), then $\mu^\omega=\mu^+$.  
For such a $\mu$, we conclude that $\mu^+\leq\mathop{cov}(\mu)\leq\mu^+=\mu^+$, giving us the value of $\mu^+$ for $\mathop{cov}(\mu)$.
By induction, we obtain the following theorem:

\begin{thm}[Gitik]\label{thm:boringcov}
    If there is no $\lambda$ with Mitchell order $\lambda^{++}$, and if $\mathfrak{c}<\aleph_\omega$, then for every uncountable cardinal $\kappa$,
    \[ \mathop{cov}(\kappa) = 
    \begin{cases}
        \kappa^+ & \text{ if }\cof(\kappa)=\omega \\
        \kappa & \text{ otherwise.}
    \end{cases}\]
\end{thm}


In the case that there is no $\kappa$ with $o(\kappa)=\kappa^{++}$, but $\mathfrak{c}>\aleph_\omega$, we have the same behavior for the $\mathop{cov}$ function starting at least as low as the least singular above $2^\omega$.





In the context of computing the sequentiality and $k$-ness numbers, our interest is in $\mathop{cov}(\kappa)$ where $\kappa$ (the size of a separable metrizable space) is no more than $\mathfrak{c}$. 
Unfortunately, then, the preceding theorem does not apply. 
\begin{qu}\label{qu:large_card}\hfill
\begin{enumerate}[noitemsep,topsep=1pt]
\item If there is no $\lambda$ with Mitchell order $\lambda^{++}$, then is $\mathop{cov}(\aleph_\omega)=\aleph_{\omega+1}$? 
\item Next we have a version of the previous item, but for other cardinals than $\aleph_\omega$ specifically: suppose that there is no $\lambda$ with Mitchell order $\lambda^{++}$, and that $\mathfrak{c}>>\aleph_{\omega}$. Is it true that for every $\omega<\kappa<\mathfrak{c}$ of countable cofinality, $\mathop{cov}(\kappa)=\kappa^+$? 
\end{enumerate}
\end{qu}

In personal communication, told to us via Cummings, Gitik thinks (\cite{Gitik:personal}) that the answer to Question A, part 2, is positive, i.e., that the \textsc{CH} is not necessary.

Assuming the existence of sufficiently large cardinals, we can construct models in which the covering number exceeds this optimal upper bound, for example at $\aleph_\omega$.
Indeed, Magidor \cite{Magidor} showed that if there is a supercompact cardinal then it is consistent that \textsc{CH} holds and $\mathop{cov}(\aleph_\omega)>\aleph_{\omega}^+$. 
Gitik 
\cite{Gitik:NegSCH} was able to lower the large cardinal strength, which combined with his result above 
gives the exact consistency strength of the failure of $\mathsf{SCH}$.

\paragraph{Controlling $\mathfrak{b}$ and $\mathfrak{c}$}
The models above where $\mathop{cov}(\kappa) > \kappa^+$ is not directly relevant for our purposes (computing $\mathop{seq}(M)$ and $k(M)$, for $M$ separable metrizable) because \textsc{CH} holds, whereas we need the continuum large while also controlling $\mathfrak{b}$. 
To correct this we modify  via forcing, intending to keep $\mathfrak{b}$ and the relevant covering numbers unchanged, while boosting the continuum.


As is well known c.c.c. forcings do not change the covering number. 
\begin{lem}\label{lem:cccPreservesCapturing}
    Let $\kappa$ be an uncountable cardinal.
        Let $\mu=\mathop{cov}(\kappa)$ and let $\mathbb{P}$ be c.c.c. Let $\dot{X}$ be the $\mathbb{P}$-name for the set of countable subsets of $\kappa$ in the extension. Then $\mathbb{P}$ forces that $\mu=\cof(\dot{X},\subseteq)$.
\end{lem}

In particular, we can use random real forcing to push the continuum high, while keeping the bounding number low, and fixing the covering number. 
\begin{thm}\label{th:small_b_big_sam}
    If there is a supercompact cardinal then it is consistent that $\mathfrak{b}=\mathfrak{d}=\aleph_1$, $\mathfrak{c}$ as large as we like, and $\mathop{cov}(\aleph_\omega)>\aleph_{\omega+1}$.
\end{thm}
\begin{proof}
    Take Magidor's model from \cite{Magidor} and recall in this model the \textsc{CH} holds, and $\mathop{cov}(\aleph_\omega)>\aleph_{\omega+1}$.
        Now do random real forcing to keep $\mathfrak{b}$ and $\mathfrak{d}$ equal to $\omega_1$, push $\mathfrak{c}$ as large as we like, and not change $\mathop{cov}(\aleph_\omega)>\aleph_{\omega+1}$. This last claim follows from Lemma~\ref{lem:cccPreservesCapturing}.
\end{proof}

A similar argument gives the following.
\begin{thm}\label{th:sam_of_b}
    If there is a supercompact cardinal then it is consistent that $\mathfrak{b}=\mathfrak{d}=\aleph_{\omega+1}$, $\mathfrak{c}$ as large as we like, and $\mathop{cov}(\aleph_{\omega+1})>\mathfrak{b}$.
\end{thm}

While from Gitik's example, Corollary~\ref{cor:noBounds}, we derive this next scenario.
\begin{thm}\label{th:sam_of_1st_fixpt}
    If there are sufficiently large cardinals then it is consistent that for $\kappa$ the first fixed point of the $\aleph$-function we have $\mathfrak{b}=\mathfrak{d}=\aleph_1$,  $\mathop{cov}(\kappa)$ as large as we like, and $\mathfrak{c}$ even larger.
\end{thm}

\paragraph{Lower Bounds of the Covering Number - Mod Finite  Scales} Below, see Section~\ref{ssec:ksam}, we will want to control the covering number not just of a single cardinal but of families. This will be achieved via lower bounds on the covering number provided by mod finite scales. 

Call a set $A$ of cardinals \emph{suitable} for a cardinal $\kappa$ of countable cofinality if $A$ is countable, every member of $A$ is uncountable, regular and $\sup A=\kappa$.
 Let $<^*$ be the mod finite order on $\prod A$, so $f<^*g$ if and only if for all but finitely many $\mu$ we have $f(\mu)<g(\mu)$.

 \begin{lem}\label{l:suitable}
     If $A$ is suitable for $\kappa$ then $([\kappa]^\omega,\subseteq) \tq (\prod A,<^*)$. 

     Hence, $\mathop{cov}(\kappa) \ge \sup \{\cof\left(\prod A,<^*\right) :A \text{ suitable for }\kappa\}$.
 \end{lem}
\begin{proof}
    For $S$ in $[\kappa]^\omega$ and $\mu$ in $A$, set $\phi_+(S)(\mu)=\sup (S \cap \mu)+1$, and note, since $S$ is countable and $\mu$ has uncountable cofinality, that this is in $\mu$; so $\phi_+$ is in $\prod A$. For $f$ in $\prod A$ set $\psi_-(f)=\mathop{im}(f)$. 
    Now if $\psi_-(f) \subseteq S$ then, for all $\mu \in A$, we have $f(\mu)$ in $S$, so $\phi_+(f)(\mu) > \sup (S \cap \mu) \ge f(\mu)$; so $f <^* \phi_+(S)$, as required.
\end{proof}

For sufficiently small $\kappa$ the above lower bound on the covering number is attained. Suppose $\kappa=\aleph_\alpha$ is such that $1 \le \alpha < \omega_1$ and $\alpha$ has cofinality $\omega$.  Then $\kappa$ has a maximum suitable set $A_\infty$, namely $A_\infty=[\aleph_1,\kappa] \cap \textsc{REG}$. By a standard \textsc{PCF} fact (see \cite{AbrahamMagidor}), $\mathop{cov}(\kappa)=\mathop{max\  pcf}(\kappa)=\cof(\prod A_\infty,<)$; this in turn equals $\cof(\prod A_\infty,<^*)$. Thus we have shown:
\begin{lem}
    If $\kappa=\aleph_\alpha$ has  countable cofinality and $1 \le \alpha < \omega_1$ then 
    \[ \mathop{cov}(\kappa) = \sup \{\cof\left(\prod A,<^*\right) :A \text{ suitable for }\kappa\}.\]
\end{lem}
 
 A (mod finite) \emph{scale} in $\prod A$ is a sequence $\vec{f}=\la f_\alpha:\alpha<\lambda\ra$ of functions in the product which is increasing mod finite ($\alpha<\beta<\lambda$ implies $f_\alpha<^*f_\beta$) and cofinal mod finite. 
 The cardinal $\lambda$ is called the length of the scale. 
 In Tukey terms, there is a scale of length $\lambda$ in $\prod A$ if and only if $\left(\prod A,<^*\right) \te \lambda$. 
If a product $\prod A$ carries a scale of length $\lambda$ regular, then its cofinality  is $\lambda$. 
 Abusing terminology, let us say that $\kappa$ admits a (mod finite) scale of length $\lambda$ if for some $A$ suitable for $\kappa$ the product $\prod A$ carries a scale of length $\lambda$. 
 By Lemma~\ref{l:suitable}, we conclude:




\begin{cor}\label{cor:ScalesIncreaseSam}
    Suppose that $\kappa$ is a singular cardinal of countable cofinality. If $\kappa$ admits a scale of regular length $\lambda$ then $\mathop{cov}(\kappa)\geq\lambda$. 
\end{cor}

 

\subsection{The Sampling Number}\label{sec:sam}


Let $\ii$ be the general relation ($\ii$ for ``infinite intersection'') defined as follows: $A  \ii   B$ if and only if $A\cap B$ is infinite. Define the \emph{sampling number} of a cardinal $\kappa$ to be $\mathop{sam}(\kappa)=\cof( [\kappa]^{\omega},\ii)$. That is, $\mathop{sam}(\kappa)$ is the least number of countable subsets of $\kappa$ needed to have infinite intersection with every countably-infinite subset of $\kappa$. 

Clearly the sampling is closely related to the covering number. 
The next result, due to Will Brian \cite{Brian}, tells us that in fact they are equal.
\begin{thm}[Brian] \label{th:cov=sam} For every infinite cardinal $\kappa$ we have 
\[([\kappa]^\omega,\ii)\te([\kappa]^\omega,\subseteq), \quad \text{ and hence }\mathop{sam}(\kappa)=\mathop{cov}(\kappa).\]
\end{thm}
\begin{proof} First note that a cofinal set for $([\kappa]^\omega,\subseteq)$ is certainly cofinal in $( [\kappa]^{\omega},\ii)$, so taking $\phi_+$ to be the identity shows $([\kappa]^\omega,\subseteq) \tq ( [\kappa]^{\omega},\ii)$. 

For the converse,  note that $\kappa$ and $[\kappa]^{<\omega}$ are bijective, so $([\kappa]^\omega,\ii)$ is isomorphic to $([[\kappa]^{<\omega}]^\omega,\ii)$. We show $([[\kappa]^{<\omega}]^\omega,\ii)\tq([\kappa]^\omega,\subseteq)$. 
    To see this, first for each $T$ in $[\kappa]^\omega$, fix an enumeration, $T=\{t_i\}_{i\in \omega}$, 
    and then  define  $\psi_-:[\kappa]^\omega \to [[\kappa]^{<\omega}]^\omega$ 
    by $\psi_-(\{t_i\}_{i \in \omega})=\{ \{t_i\}_{i  \le n}\} : n \in \omega\}$, and $\phi_+:[[\kappa]^{<\omega}]^\omega \to [\kappa]^\omega$ by $\phi_+(S)=\bigcup S$. 
    Now if $\psi_-(\{t_i\}_{i \in \omega}) \ii S$, then for infinitely many $n$, $\{t_i\}_{i \le n} \in S$, so $t_0,\ldots,t_n \in \bigcup S=\phi_+(S)$, and hence, $\{t_i\}_{i \in \omega} \subseteq \phi_+(S)$, as required. 
\end{proof}
    






 












\section{The Minimum Size of Generating Families} 
\subsection{How Many Sequences Suffice?}\label{ssec:seqsam}
\paragraph{The $\mathbf{seq}(M)$ Relation} 
Let $M$ be a (separable metrizable) space, and let $\nci$ be the general relation defined for subsets $S$ and $T$ of $M$ by  
$S \nci T$ if and only if $S \cap T$ is not closed in $T$. 
Set $\mathop{NC}(M)$ to be all subsets of $M$ that are not closed, and let   
 $\mathop{CS^+}(M)$ be all infinite convergent sequences with limit. 
Define $\mathbf{seq}(M)=(NC(M),\mathop{CS}^+(M),\nci)$,
the \emph{sequential structure} of $M$, and  note that its cofinal sets correspond to generating collections of (infinite) convergent sequences. 
Hence our invariant $\mathop{seq}(M)$ is the cofinality of the relation $\mathbf{seq}(M)$. From \cite{FG_Shape} we have an explicit, combinatorial description of $\mathbf{seq}(M)$ up to Tukey-equivalence which is summarized in the next paragraph. Below we use this to compute $\mathop{seq}(M)$. 
We will find that $\mathop{seq}(M)$ is ranked by the size of $M$: if $|M| \le |N|$ then $\mathop{seq}(M)\le\mathop{seq}(N)$. Consequently we will compute $\mathop{seq}(M)$ first for countable spaces, before moving on to the uncountable case.



\paragraph{Characterization and Realization of $\mathbf{seq}(M)$}
For a point $x$ in $M$ fix $U_x$ an open neighborhood of $x$ of minimal cardinality. 
 Let $M_\infty$  be the set of $x$ in $M$ with $|U_x|$ equal to $|M|$. 
Separate the points of $M_\infty$ in two as follows. Let $M_\infty^+$ be all $x$ from $M_\infty$ which have a neighborhood base $(B_n)_n$ such that every $D_n=B_n \setminus B_{n+1}$ has cardinality $|M|$. Let $M_\infty^-$ be $M_\infty \setminus M_\infty^+$.
Then a point $x$ from $M_\infty$ is in $M_\infty^-$ if and only if it has a neighborhood base $(B_n)_n$ such that every $D_n=B_n \setminus B_{n+1}$ has cardinality $<|M|$. 
Let $\lambda_M=|M_\infty^+|$ and $\mu_M=|M_\infty^-|$. 
The cardinals $|M|,\lambda_M$ and $\mu_M$ suffice to characterize $\mathbf{seq}(M)$ up to Tukey equivalence.


\begin{thm}\label{thm:seqTukeyEquivalence} (\cite{FG_Shape}, Theorem~4.10) Let $M$ be separable metrizable.  Then $\lambda_M \le |M|$, $\mu_M \le \omega$ and $\mathbf{seq}(M)$ is Tukey equivalent to 
    \[
    P(|M|) \times (\lambda_M,=) + Q(|M|) \times (\mu_M,=) + S(|M|).
    \]   
\end{thm}

The case when $\lambda_M$ and $\mu_M$ are both equal to zero is an important special case. 
Define a space $M$ to be \emph{locally small} if  every point of $M$ has a neighborhood of size strictly less than $|M|$, and note  
this is equivalent to saying $\lambda_M=0=\mu_M$. 
Note as well that since the neighborhood may be assumed to come from any given base, and separable metrizable spaces have  countable bases, it follows that if $M$ is a locally small, separable metrizable space then $|M|$ has countable cofinality. 

Every Tukey type mentioned in the characterization can be realized as the $\mathbf{seq}(M)$ of a separable metrizable space. 
\begin{thm}\label{th:seq_real}(\cite{FG_Shape}, Theorem~4.11) 

For every cardinal $\kappa \le \mathfrak{c}$ there is a separable metrizable space $M=M(\kappa)$ such that $\mathbf{seq}(M) \gte P(\kappa) \times (\kappa,=)$. 

    For every cardinal $\kappa \le \mathfrak{c}$ of countable cofinality and $\lambda < \kappa$ and $\mu \in \omega+1$ there is a separable metrizable space $N=N(\kappa,\lambda,\mu)$ such that $\mathbf{seq}(N) \gte P(\kappa) \times (\lambda,=) + Q(\kappa) \times (\mu,=) + S(\kappa)$. 
\end{thm}



\paragraph{The Countable Case - Countable Spaces}

Let $M$ be a countably infinite, separable metrizable space. In this case the relation $S(|M|)=S(\omega)$ is trivial, because $P(n)$ is trivial for every $n$ in $\omega$. Recall $\cof(P(\omega))=\mathfrak{b}$. 
Further, $M$ partitions into three subsets: (1) $M\setminus M_\infty$, which are the isolated points of $M$, 
(2) $M_\infty^-$, which are the points of $M$ with a neighborhood homeomorphic to the convergent sequence, 
and (3) the remaining points, $M_\infty^+$. 
Putting the data together we compute $\mathop{seq}(M)$ in the countably infinite case.

\begin{thm}\label{th:seq_ctble} Let $M$ be a countably infinite, separable metrizable space. Then
\[
\mathop{seq}(M) = 
\begin{cases}
  0 & \text{ if $M$ is discrete} \\
  n & \text{ if $M$ is the disjoint sum of $n$ $(\le \omega)$ convergent sequences} \\
  \mathfrak{b} & \text{ otherwise.}
\end{cases}
\]    
\end{thm}

\paragraph{The Uncountable Case - Connection to the Sampling Number}

It remains to consider uncountable separable metrizable $M$. The first step is to compute the cofinalities of $P(\kappa)$ and $Q(\kappa)$, for some uncountable cardinal $\kappa$ (soon to be $|M|$), in terms of the sampling number of $\kappa$, $\mathop{sam}(\kappa)$, and the bounding number, $\mathfrak{b}$,
\begin{thm}\label{th:cof_Pkappa}
Let $\kappa$ be an uncountable cardinal, and $(\kappa_n)_n$ a sequence of infinite cardinals increasing (not necessarily strictly) to $\kappa$. 
Then 
\[\cof((\prod_n [\kappa_n]^{<\omega})_\infty,i_\infty)
= \mathop{sam}(\kappa)\cdot \mathfrak{b}. 
\]
Hence $\cof(P(\kappa))= \mathop{sam}(\kappa)\cdot\mathfrak{b}$, and if $\cof(\kappa)=\omega$, then $\cof(Q(\kappa))=\mathop{sam}(\kappa)\cdot\mathfrak{b}$.    
\end{thm}



\begin{proof}
It suffices to show the following pair of inequalities: 

(1)~$\cof(\left(\prod_n [\kappa_n]^{<\omega}\right)_\infty,i_\infty) \ge \mathop{sam}(\kappa)\cdot\mathfrak{b}$, and 
(2)~$\mathop{sam}(\kappa)\cdot\mathfrak{b}  \ge \cof(P(\kappa))$.

\smallskip

For (1), first note that each $\kappa_n \ge \omega$, so $(\left(\prod_n [\kappa_n]^{<\omega}\right)_\infty,i_\infty) \gtq P(\omega) \gte (\omega^\omega,\le_\infty)$, and the last relation has cofinality $\mathfrak{b}$. 
We get (1) by combining this with the fact that  $(\left(\prod_n [\kappa_n]^{<\omega}\right)_\infty,i_\infty) \gtq ([\kappa]^\omega,\ii)$. 
To see the latter, define $\phi_+( (F_n)_n)$ to be $\bigcup_n F_n$ when this set is infinite, and an arbitrary element of $[\kappa]^\omega$ otherwise. 
For any $S$ in $[\kappa]^\omega$, enumerate $S$ as $S=\{s_n : n \in \omega\}$, and set $\psi_-(S)=( \{s_n\})_n$. 
Now if $\psi_-(S) i_\infty (F_n)_n$, then for infinitely many $n$ we have $s_n \in F_n$; it follows $\bigcup_n F_n$ is infinite and meets $S$ in an infinite set, so $S \ii \phi_+((F_n)_n)$, as required.

Now we establish (2). First fix $\mathcal{S}$ cofinal in $([\kappa]^\omega,\ii)$ of size $\mathop{sam}(\kappa)$. 
Take any $S$ in $\mathcal{S}$, and enumerate it as $S=\{s_n : n \in \omega\}$. Observe that the bijection $n \leftrightarrow s_n$ induces a Tukey equivalence between $( ([S]^{<\omega})^{\omega}_{\infty}, i_\infty)$ and $( ([\omega]^{<\omega})^{\omega}_{\infty}, i_\infty)$; and so both have cofinality $\mathfrak{b}$. Fix $\mathcal{C}_S$ cofinal in $( ([S]^{<\omega})^{\omega}_{\infty}, i_\infty)$ of size $\mathfrak{b}$. 

Let $\mathcal{C}= \bigcup \{ \mathcal{C}_S : S \in \mathcal{S}\} \cup \mathcal{C}_0$, where $\mathcal{C}_0 = \{ (F)_n : \emptyset \ne F \in [\kappa]^{<\omega}\}$. 
Note that $\mathcal{C}$ has size $\mathfrak{b} \cdot \cof( [\kappa]^{\omega},\ii)\cdot\kappa$. 
We show $\mathcal{C}$ is cofinal in $P(\kappa)$. Since $\cof(P(\kappa))$ is at least $\kappa$, by Lemma \ref{lem:cofatleastkappa}, 
Claim (2) then follows. 


Take any $(G_n)_n$ in $([\kappa]^{<\omega})^{\omega}_{\infty}$. Let $T=\bigcup_n G_n$. If $T$ is finite, then $(T)_n$ is in $\mathcal{C}_0$ and $(G_n)_n i_\infty (T)_n$. 
So suppose $T$ is infinite. Pick $S$ in the cofinal set $\mathcal{S}$ so that $T'=S \cap T$ is infinite. 
For each $n$ set $G'_n=G_n \cap T'$. Since $T'$ is infinite, for infinitely many $n$ we must have $G'_n$ non-empty, so  $(G'_n)_n$ is in $([S]^{<\omega})^{\omega}_{\infty}$.  
As $\mathcal{C}_S$ is cofinal in $(([S]^{<\omega})^{\omega}_{\infty},i_\infty)$, there is a $(F_n)_n$ in $\mathcal{C}_S$ so that $(G'_n)_n i_\infty (F_n)_n$, and so $(G_n)_n i_\infty (F_n)_n$,
as required for $\mathcal{C}$ to be cofinal.
\end{proof}

Combining the previous theorem with the Tukey characterization of $\mathbf{seq}(M)$, Theorem~\ref{thm:seqTukeyEquivalence} above, we compute $\mathop{seq}(M)$, for uncountable separable metrizable $M$, in terms of the /covering number of the size of $M$ and $\mathfrak{b}$. 
\begin{thm}\label{th:seq_sam}
     Let $M$ be an uncountable separable metrizable space. 
\begin{enumerate}[noitemsep,topsep=1pt]
\item[(i)] If $M$ is not locally small then 
   $\mathop{seq}(M)= \mathop{cov}(|M|) \cdot \mathfrak{b}$.
     
\item[(ii)]      If $M$ is locally small then $\mathop{seq}(M)=\lim_{\mu <|M|} \mathop{cov}(\mu)\cdot \mathfrak{b}$.    
\end{enumerate}
   \end{thm}

   \begin{proof}
   Let $\kappa=|M|$.  If it has countable cofinality, then fix a strictly increasing sequence of infinite cardinals, $(\kappa_n)_n$, with limit $\kappa$. Otherwise let $\kappa_n=\kappa$ for all $n$, and note $(\left(\prod_n [\kappa_n]^{<\omega}\right)_\infty,i_\infty)= P(\kappa)$.
From Theorem \ref{thm:seqTukeyEquivalence}, we know $\mathbf{seq}(M)$ is Tukey equivalent to 
$P(\kappa) \times (\lambda_M,=) + Q(\kappa) \times (\mu_M,=) + S(\kappa)$,  
where $\lambda_M \le \kappa$ and $\mu_M \le \omega$. Also recall that when $\kappa$ has uncountable cofinality, $Q(\kappa)$ and $S(\kappa)$ are trivial.

If $\kappa$ has countable cofinality then there are two cases. In the first, $\lambda_M=\mu_M=0$. Again by Theorem \ref{thm:seqTukeyEquivalence}, this occurs when $M$ is locally small, in which case $\mathbf{seq}(M)$ is Tukey equivalent to $S(\kappa)=\sum_n P(\kappa_n) \times (\kappa_n,=)$. Now 
\begin{multline*}
\mathop{seq}(M)=\cof(S(\kappa))=\lim_n \cof(P(\kappa_n))\cdot\kappa_n =\lim_n\cof(P(\kappa_n)) \\
=\lim_n\mathop{sam}(\kappa_n)\cdot\mathfrak{b}=\lim_{\mu <|M|} \mathop{sam}(\mu)\cdot \mathfrak{b};
\end{multline*}
note that the second equality uses the definition of $S(\kappa)$, that the third uses Lemma~\ref{lem:cofatleastkappa}, that the next equality follows from Theorem~\ref{th:cof_Pkappa}, and that the final equality uses the fact that the $\kappa_n$ are cofinal in $\kappa$ and that the $\mathop{sam}$ operator is monotone (Lemma~\ref{lem:samIncreasing}).

If $\kappa$ has countable cofinality but at least one of $\lambda_M, \mu_M \ne 0$, or when $\kappa$ has uncountable cofinality, then $P(\kappa)\times (\kappa,=) \gtq \mathbf{seq}(M) \gtq (\left(\prod_n [\kappa_n]^{<\omega}\right)_\infty,i_\infty)$, and so the claim follows from Theorem~\ref{th:cof_Pkappa}.    
\end{proof}
   

%\subsection{Sampling In ZFC and When `No Large Cardinals'}





%If there are `no large cardinals' then we have equality, $\mathop{sam}(\aleph_\omega) = \aleph_{\omega+1}$. Then, by claims (i) and (iii), we see, for every $n \in \mathbb{N}$ we have  $\mathop{sam}(\aleph_{\omega+n}) = \aleph_{\omega+n}$. 
%With a sufficiently strong `no large cardinals' hypothesis this pattern repeats, for every uncountable cardinal $\kappa$ we have $\mathop{sam}(\kappa) = \kappa$, unless $\cof(\kappa)=\omega$, in which case $\mathop{sam}(\kappa) = \kappa^+$. 
%The above claims follow from Shelah's \textsc{PCF} theory.

 


%\subsection{Pushing Up Sampling Numbers}




  



\paragraph{Computing $\mathop{seq}(M)$}

For $\mathop{seq}(M)$ of countable separable metrizable $M$, see Theorem~\ref{th:seq_ctble}. Here we summarize our understanding of $\mathop{seq}(M)$ for uncountable $M$.
\begin{itemize}[topsep=2pt]
\item We can compute $\mathop{seq}(M)$ from $\mathop{cov}(|M|)$ and $\mathfrak{b}$: 
 $\mathop{seq}(M)=\mathop{cov}(|M|)\cdot\mathfrak{b}$, unless $M$ is locally small when  
$\mathop{seq}(M)=\lim_{\mu <|M|} \mathop{cov}(\mu)\cdot \mathfrak{b}$.  



\item All possible values of $\mathop{seq}(M)$ are attained: if $\kappa \le \mathfrak{c}$ is uncountable, then there is an $M(\kappa)$ with  $\mathop{seq}(M(\kappa))=\mathop{cov}(\kappa)\cdot \mathfrak{b}$; if in addition, $\kappa$ has countable cofinality, then there is an $N(\kappa)$ with $\mathop{seq}(N(\kappa))=\lim_{\mu <\kappa} \mathop{cov}(\mu)\cdot \mathfrak{b}$. This follows \textsl{a fortiori} from Theorem~\ref{th:seq_real}. 

\item The sequentiality number is ranked by the size of the space: if $|M| \le |N|$ then $\mathop{seq}(M) \le \mathop{seq}(N)$. Further, if $\kappa \le \mathfrak{c}$ is uncountable, then there is exactly one value of $\mathop{seq}(M)$ for $M$ of size $\kappa$, unless $\cof(\kappa)=\omega$, in which case at least one but possibly two values attained (one for locally small $M$, and a potentially strictly larger value for all other $M$ of size $\kappa$).   
\end{itemize}

To go deeper, observe that two cases arise depending on the relative sizes of $\mathop{cov}(|M|)$ and $\mathfrak{b}$.
 If $\mathop{cov}(|M|) \le \mathfrak{b}$ then $\mathop{seq}(M)=\mathfrak{b}$. 
In particular, if $\mathfrak{b}=\mathfrak{c}$ then $\mathop{seq}(N)=\mathfrak{c}$ for all uncountable separable metrizable $N$. 

Now suppose $\mathfrak{b} \le \mathop{cov}(|M|)$. In this case the following hold.
\begin{itemize}[topsep=2pt]
\item Always $\mathop{seq}(M) \ge |M|$, and   if $|M|=\aleph_n$ or $|M|=\mathfrak{c}$ then $\mathop{seq}(M)=|M|$.

\item If $\cof(|M|)=\omega$ but $M$ is not locally small, then $\mathop{seq}(M) \ge |M|^+$.

\item If $|M|=\aleph_\alpha$ is strictly below the first fixed point of the $\aleph$-function then $\mathop{seq}(M) < \aleph_{|\alpha|^{+4}}$. 

\item It is consistent, assuming the existence of sufficiently large cardinals, that there is a separable metrizable $M$ of size the first fixed point of the $\aleph$-function, with $\mathop{seq}(M)$ as far above $|M|$ as we like. 

\item It is consistent, assuming the existence of a supercompact cardinal, that there is a separable metrizable $M$ with $|M|=\aleph_\omega$ but $\mathop{seq}(M) > |M|^+$. 

\item It is consistent, assuming the existence of a supercompact cardinal, that there is a separable metrizable $M$ with $|M| = \mathfrak{b}$ but $\mathop{seq}(M) > \mathfrak{b}$. 
\end{itemize}

It is an open problem whether large cardinals are necessary for the last two points. Is it true, under some `no large cardinals' hypothesis, that $\mathop{seq}(M)=|M|$ if $\cof(|M|)>\omega$ or $M$ is locally small, and otherwise $\mathop{seq}(M)=|M|^+$?
See Question~\ref{qu:large_card}. 





\subsection{How Many Compact Sets Suffice?}\label{ssec:ksam}

\paragraph{The $\mathbf{k}(M)$ Relation}
The \emph{$k$-structure} of a (separable metrizable) space $M$ is the relation $\mathbf{k}(M)=(\mathop{NC}(M),\K(M),\nci)$, where $\mathop{NC}(M)$  are the non-closed subsets of $M$, the general relation $\nci$ is ``not closed in'' and $\K(M)$ is the collection of all compact subsets of $M$.  
Observe  that cofinal sets of $\mathbf{k}(M)$ correspond to generating compact collections, and hence $\mathop{k}(M)=\cof(\mathbf{k}(M))$. Clearly, $\mathbf{seq}(M) \tq \mathbf{k}(M)$, and so $\mathop{seq}(M) \ge \mathop{k}(M)$. 
Additionally, define the \emph{$kc$-structure}  $\mathbf{kc}(M)=(M,\K(M),\in)$, whose  cofinal sets are the compact covers of $M$,  and the \emph{compact-covering number}, $\mathop{kc}(M)=\cof(\mathbf{kc}(M))$, the minimum number of compact sets needed to cover $M$ (see \cite{FG_Shape_CC}). Note $kc(M)\le \omega$ if and only if $M$ is $\sigma$-compact.
From Lemma~3.5 of \cite{FG_Shape} we know that $\mathbf{k}(M) \tq \mathbf{kc}(M)$, and so $\mathop{k}(M) \ge kc(M)$, for every separable metrizable space, $M$, which is not $\sigma$-compact. 

Our results above on $\mathop{seq}(M)$ -  building on the characterization of $\mathbf{seq}(M)$ up to Tukey equivalence obtained in \cite{FG_Shape}, along with what is known about the covering number - tell us the following: that $\mathop{seq}(M)$ is ranked by the size of $M$, if $|M|\le|N|$ then $\mathop{seq}(M)\le\mathop{seq}(N)$, and for a fixed size of space at most two values of $\mathop{seq}(M)$ can occur:  \ for every uncountable $\kappa \le \mathfrak{c}$
\[ \{\mathop{seq}(M) : |M|=\kappa\} \subseteq \left\{\lim_{\mu <|M|} \mathop{cov}(\mu) \cdot\mathfrak{b},\mathop{cov}(|M|)\cdot\mathfrak{b}\right\}.\] 



The results of \cite{FG_Shape} give in complete detail the Tukey types of $\mathbf{k}(M)$ for $M$ $\sigma$-compact, from which we deduce the corresponding values of $k(M)$ in the next paragraph. We turn then to the case when $kc(M)$ is uncountable. 
However simple descriptions of the Tukey type of $\mathbf{k}(M)$ are not known for spaces which are not $\sigma$-compact. To the contrary, from \cite{FG_Shape} we know that the structure of Tukey types of $\mathbf{k}(M)$'s is immensely complex. There is a Tukey order anti-chain of size $2^\mathfrak{c}$, and both $\mathfrak{c}^+$ and the unit interval (with the usual order) order embed in $\mathbf{k}(M)$'s under the Tukey order.

Nevertheless we show below that $kc(M)$ plays a role for $k(M)$ loosely analogous to the role that $|M|$ plays for $\mathop{seq}(M)$, with the connection being made by the covering number combined with the bounding number. 
Theorem~\ref{th:k_sam} below implies that:  \ for every uncountable $\kappa \le \mathfrak{c}$
\[\{\mathop{k}(M) : kc(M)=\kappa\} \subseteq [\kappa\cdot\mathfrak{b},\mathop{cov}(\kappa)\cdot\mathfrak{b}].\]


Above we explored, in the context of $\mathop{seq}(M)$, the possible values of $\mathop{cov}(\kappa)$ and the gap between $\kappa$ and $\mathop{cov}(\kappa)$, for $\kappa=|M|$. We do not rehash here the situation when $\kappa=kc(M)$. Instead in the sequel we focus on what is different between  $\mathop{seq}(M)$ and $k(M)$, 
namely (1)  that while $\mathop{seq}(M)$  for $|M|=\kappa$ can only take on the value of one of a lower and an upper bound,  
$\{\mathop{k}(M) : kc(M)=\kappa\}$ is constrained merely to an \emph{interval},  and (2) the lower bound ($\kappa\cdot\mathfrak{b}$) on this interval of $k(M)$ values is potentially smaller than the analogous lower bound value for $\mathop{seq}(M)$ values ($\lim_{\mu <\kappa} \mathop{cov}(\mu)\cdot\mathfrak{b}$). 
Our task is made more difficult because we have no analog to the realization of types theorem, Theorem~\ref{th:seq_real}, for $\mathbf{seq}(M)$. 

Let $\kappa$ be an uncountable cardinal  with $\kappa \le \mathfrak{c}$.  Define the \emph{$k$-spread} at $\kappa$, to be $\sigma_k(\kappa)=|\{k(M) : kc(M)=\kappa\}|$. We show in Theorem~\ref{th:k=sam} that 
the upper bound of the interval, $[\kappa\cdot\mathfrak{b},\mathop{cov}(\kappa)\cdot\mathfrak{b}]$, is always attained; that is, for some $M$ we have $kc(M)=\kappa$ and $k(M)=\mathop{cov}(\kappa)\cdot\mathfrak{b}$, and hence $\sigma_k(\kappa) \ge 1$. 
We then show, Theorem~\ref{th:powers}, that in certain circumstances the lower bound, $\kappa\cdot\mathfrak{b}$, is also attained. It follows that $\sigma_k(\kappa)$ may be at least two. 
Finally we show below that, modulo large cardinals, for every $2 \le K < \omega$ it is consistent that for some $\kappa$ we have 
$\sigma_k(\kappa)=K$ (Theorem~\ref{th:K_many})  and it is consistent that for some $\kappa$ we have 
$\sigma_k(\kappa)\ge \aleph_0$ (Theorem~\ref{th:ctbl_many}). 



 

\paragraph{The Countable Case - $\sigma$-Compact Spaces}

Van Douwen computed $\mathop{k}(M)$ for $\sigma$-compact separable metrizable $M$ (and this data can also be read off from the classification of Tukey types of $\mathbf{k}(M)$ in \cite{FG_Shape} for these $M$). 

\begin{thm}
    Let $M$ be $\sigma$-compact and separable metrizable. Then
\[
\mathop{k}(M) = 
\begin{cases}
  0 & \text{ if $M$ is discrete} \\
  1 & \text{ if $M$ is compact} \\
  \omega & \text{ if $M$ is locally compact but not compact} \\ 
  \mathfrak{b} & \text{ otherwise.}
\end{cases}
\]    
\end{thm}

\paragraph{Another Connection To Sampling}

We now see that for non $\sigma$-compact separable metrizable spaces the $k$-ness number is connected to the compact-covering number via the covering number.

\begin{thm}\label{th:k_sam}
   For every separable metrizable $M$ which is not $\sigma$-compact 
   \[\mathop{cov}(kc(M))\cdot\mathfrak{b} \ge k(M) \ge kc(M)\cdot\mathfrak{b}.\]   
\end{thm}
\begin{proof} Since $M$ is not $\sigma$-compact we know, by Lemma~3.5 and Theorem~3.11 of \cite{FG_Shape}, $\mathbf{k}(M) \gtq \mathbf{kc}(M)$ and $\mathbf{k}(M) >_T (\omom,\le_\infty)$, so $k(M) \ge kc(M)\cdot\mathfrak{b}$.

Now let $\kappa=kc(M)$, and fix  $\mathcal{K}=\{K_\alpha : \alpha \in \kappa\}$  a compact cover of $M$. We will show that $(\kappa,=) \times (\left( [\kappa]^{<\omega} \right)^{\omega}_{\infty},i_\infty) \tq \mathbf{k}(M)$. 
Then $\mathop{sam}(kc(M))\cdot\mathfrak{b} \ge k(M)$ follows from 
Theorem~\ref{th:cof_Pkappa}.

    For each $x$ in $M$, fix $\alpha(x)$ in $\kappa$ so that $x \in K_{\alpha(x)}$. For any finite subset $F$ of $\kappa$, let $K_F$ be the compact set $K_F:=\bigcup_{\alpha \in F} K_\alpha$. 
Fix a metrizable compactification $\gamma M$ of $M$, with compatible metric $d$, and select a base $\mathcal{B}=\{B_m : m \ge 1\}$ for $\gamma M$ where $B_1=\gamma M$, $\mathop{diam} B_m \to 0$ and for every $x$ in $\gamma M$ the set $\{n : x \in B_n\}$ is infinite.

    Define $\phi_+(\alpha, (F_n)_n)$ to be $K_\alpha \cup \bigcup \{ K_{F_m} \cap \cl{B_m} : m \ge 1 \& K_\alpha \cap \cl{B_m} \ne \emptyset\}$. 
We check  $\phi_+(\alpha,(F_n)_n)$ is compact. 
Take any cover by basic open sets. Some finite subcollection, say $\mathcal{U}_1$, covers the compact set $K_\alpha$. 
As $K_\alpha$ and $\gamma M \setminus \bigcup \mathcal{U}_1$ are compact and disjoint, the distance between them is strictly positive. 
As $\mathop{diam} B_m \to 0$, it follows that,  of the $\cl{B_m}$ that meet $K_\alpha$, all but finitely many  are contained in $\bigcup \mathcal{U}_1$. 
Hence $\phi_+(\alpha,(F_n)_n) \setminus \bigcup  \mathcal{U}_1$ is compact, and can be covered with a finite subcollection, say $\mathcal{U}_2$, from the basic open cover. Now $\mathcal{U}_1 \cup \mathcal{U}_2$ is a finite subcover of $\phi_+(\alpha,(F_n)_n)$.
    
Having defined $\phi_+$, we now work to define $\psi_-$. Let $N$ be in $NC(M)$. 
    Pick $S=(x_n)_n$ an infinite sequence of points from $N$ converging, say, to $x\notin N$. Let $\alpha=\alpha(x)$. Extract a subsequence from the  base giving a local base at $x$, $(B_{m_n})_n$, such that:  $(m_n)_n$ is strictly increasing, $m_1=1$,  $\cl{B_{m_{n+1}}} \subseteq B_{m_n}$ for all $n$, and $S \cap (\cl{B_{m_n}} \setminus B_{m_{n+1}})$ is non-empty for all $n$. Note also that $S \cap (\cl{B_{m_n}} \setminus B_{m_{n+1}})$ is finite, since otherwise $\cl{B_{m_n}} \setminus B_{m_{n+1}}$ would contain $x$, the limit of $S$. 
    Set $G_{m_n}=\{\alpha(x') : x' \in S \cap (\cl{B_{m_n}} \setminus B_{m_{n+1}})\}$, and by the preceding comments,  this set is finite and non-empty. 
    For all other $m$ set $G_m=\emptyset$. Then define $\psi_-(N) = (\alpha,(G_m)_m)$.

    We verify that $(\psi_-,\phi_+)$ are a Tukey morphism pair. 
    Suppose that 
    $
    \psi_-(N) (=\times\ii) (\alpha,(F_m)_m)$. 
    Then $\alpha=\alpha(x)$, so $x$ is in $K_\alpha$; further for infinitely many $m$, $G_m$ meets $F_m$. 
    Since $G_m=\emptyset$ when $m$ is not an $m_n$, it follows, for infinitely many $n$, $G_{m_n}$ meets $F_{m_n}$. Thus, looking at the definition of $\phi_+(\alpha,(F_m)_m)$, and choice of $(B_{m_n})_n$, we see $S \cap \phi_+(\alpha,(F_m)_m)$ is infinite, and so not closed in $\phi_+(\alpha,(F_m)_m)$. \textsl{A fortiori}, $N \cap \phi_+(\alpha,(F_m)_m)$ is not closed in $\phi_+(\alpha,(F_m)_m)$, as required.
\end{proof}






\paragraph{Totally Imperfect - The Upper Bound is Always Attained}

A separable metrizable space $M$ is \emph{totally imperfect} if every compact subset of $M$ is countable, equivalently - since uncountable compact metrizable spaces contain a Cantor set - if $M$ contains no  compact subsets of size continuum. 
Note that if $M$ has size strictly less than the continuum then $M$ is totally imperfect; in particular, locally small spaces are totally imperfect (if $M$ is locally small then $\cof(|M|)=\omega$, but $\cof(\mathfrak{c})\ne\omega$). 
However there are totally imperfect separable metrizable spaces of size continuum, Bernstein subsets of the reals, for example. 
We can determine the $k$-ness number of totally imperfect spaces exactly in terms of $\mathop{cov}(kc(M))$ and the bounding number.

\begin{thm}\label{th:k=sam}
Let $M$ be a non $\sigma$-compact, separable metrizable space.
\begin{enumerate}[noitemsep,topsep=1pt]
\item[(i)] If $M$ is totally imperfect but not locally small then 
$k(M)=\mathop{cov}(kc(M))\cdot\mathfrak{b}$.

\item[(ii)] If $M$ is locally small then 
$
k(M)=\lim_{\mu < kc(M)} \mathop{cov}(\mu)\cdot\mathfrak{b}$. 
\end{enumerate}
\end{thm}
\begin{proof}
We begin with item (i). Since $M$ is not $\sigma$-compact, Theorem \ref{th:k_sam} implies that $\mathop{sam}(kc(M))\cdot\mathfrak{b} \ge k(M) \ge \mathfrak{b}$.
As $M$ is totally imperfect, all compact subsets of $M$ are countable, and hence $kc(M)=|M|$ (i.e., we need $|M|$-many compact sets to cover $M$). As $M$ is not locally small, by Theorem~\ref{th:seq_sam}, we have $\mathop{seq}(M)=\mathop{sam}(|M|)\cdot\mathfrak{b}$, which from what we just observed, in turn equals $\mathop{sam}(kc(M))\cdot\mathfrak{b}$. 
So it suffices to show that $k(M)\cdot\mathfrak{b} \ge \mathop{seq}(M)$.

Let $\K$ be a family of compact subsets of $M$ which has size $k(M)$ and is generating. Each $K$ in $\K$ is countable, and so by Theorem~\ref{th:seq_ctble}, $\mathop{seq}(K) \le \mathfrak{b}$. Accordingly, fix $\mathcal{S}_K$ which has size no more than $\mathfrak{b}$ and is cofinal in $\mathbf{seq}(K)$. 
We claim that $\mathcal{S}=\bigcup \{ \mathcal{S}_K : K \in \K\}$, which is a subset of $CS^+(M)$ of size no more than $k(M)\cdot\mathfrak{b}$, is cofinal in $\mathbf{seq}(M)$. Indeed, take any $N$ in $NC(M)$. There is a $K$ in $\K$ so that $N'=N \cap K$ is not closed in $K$. Then there is an $S$ in $\mathcal{S}_K$ with $N' \cap S$ not closed in $S$. Now $S$ is in $\mathcal{S}$ and $N\cap S =N'\cap S$ is not closed in $S$, as required.

\smallskip

Now for (ii). Let $M$ be locally small. Then, by definition of `locally small', the Lindel\"{o}f property and the regularity of $M$, there is a countable open cover, $\mathcal{U}=\mathcal{U}_M$, such that for every $U$ from $\mathcal{U}$, we have $|\cl{U}| < |M|$. Since $\cal{U}$ is countable, we then obtain 
$
|M|=(\sup_{} \{|\cl{U}|: U \in \mathcal{U}\})\cdot\aleph_0=\{|\cl{U}|:U \in \mathcal{U}\}$,
and hence $\cof(|M|)=\omega$. Since $M$ is a separable metrizable space, $|M|\leq\mathfrak{c}$, and as $\mathfrak{c}$ has uncountable cofinality, we conclude that $|M|<\mathfrak{c}$. It then follows that $M$ is totally imperfect; as in the first paragraph, this in turn implies that $kc(M)=|M|$. 
We need to show, for such a space, both of (a)${}_M$  $k(M) \ge \lim_{\mu < |M|} \mathop{sam}(\mu) \cdot\mathfrak{b}$ and (b)${}_M$ $k(M) \le \lim_{\mu < |M|} \mathop{sam}(\mu)\cdot\mathfrak{b}$.

Towards (a)${}_M$, suppose, inductively that $M$ is locally small, and that for every locally small $M'$ of strictly smaller size, (a)${}_{M'}$ holds. 
Take any $\mu < |M|$, and we will show $k(M) \ge \mathop{sam}(\mu)\cdot\mathfrak{b}$. 
Let $\mathcal{U}=\mathcal{U}_M$, as above, and pick $U$ in $\mathcal{U}$ such that $\mu < |\cl{U}| < |M|$. 
Set $M'=\cl{U}$. As $M'$ is a closed subset of $M$, $k(M) \ge k(M')$. 
If $M'$ is not locally small, then by part (i), $k(M') = \mathop{sam}(|M'|)\cdot\mathfrak{b} \ge \mathop{sam}(\mu)\cdot\mathfrak{b}$, as required. 
Otherwise $M'$ is locally small, and then by the inductive hypothesis, $k(M') \ge \lim_{\mu' < |M'|} \mathop{sam}(\mu')\cdot\mathfrak{b} \ge \mathop{sam}(\mu)\cdot\mathfrak{b}$ (as $\mu < |M'|$), again as required for (a)${}_M$.

To complete the proof of (ii), we verify (b)${}_M$. 
Let $\mathcal{U}=\mathcal{U}_M$, as above. For each $U$ from $\mathcal{U}$, $\cl{U}$ is totally imperfect, since $M$ is. Once more arguing as in the first paragraph, we conclude that $kc(\cl{U})=|\cl{U}|$. We thus know by Theorem~\ref{th:k_sam} that $k(\cl{U}) \le \mathop{sam}(|\cl{U}|)\cdot\mathfrak{b}$, so fix $\mathcal{C}_U$ cofinal in $\mathbf{k}(\cl{U})$ of size no more than $\mathop{sam}(|\cl{U}|)\cdot\mathfrak{b}$. 
Set $\mathcal{C}=\bigcup \{ \mathcal{C}_U : U \in \mathcal{U}\}$, and note that $|\mathcal{C}| \le \lim_{\mu < |M|} \mathop{sam}(\mu) \cdot\mathfrak{b}$. We check $\mathcal{C}$ is cofinal in $\mathbf{k}(M)$, and so confirm (b)${}_M$. 
Take any $N$ in $NC(M)$. Pick a sequence $S$ on $N$ converging to a point $x$ outside $N$. Then $x$ is in some $U$ from the cover $\mathcal{U}$, and - dropping finitely many terms if needed - we can suppose $S$ is contained in $U$. Now some member of $\mathcal{C}_U$, and so $\mathcal{C}$,  meets $S$ infinitely often, as required for cofinality.
\end{proof}

 
It follows from (i) that the top of the range, $\mathop{cov}(\kappa)\cdot\mathfrak{b}$, is always obtained, in \textsc{ZFC}.
\begin{cor}
    If $\omega < \kappa \le \mathfrak{c}$, then there is a separable metrizable space $M_\kappa$ so that $k(M_\kappa)=\mathop{cov}(\kappa)\cdot\mathfrak{b}$.
\end{cor}
To see this, if $\kappa<\mathfrak{c}$ then let $M_\kappa$ be of size $\kappa$, but not locally small, and otherwise  let $M_\mathfrak{c}$ be totally imperfect of size $\mathfrak{c}$ (a Bernstein set, for example).

Similarly, from (ii) (taking $M_\kappa$ to be locally small of size $\kappa$) we have:
\begin{cor}
    If $\omega < \kappa \le \mathfrak{c}$ and $\cof(\kappa)=\omega$, then there is a separable metrizable space $M_\kappa$ so that $k(M_\kappa)=\lim_{\mu < \kappa} \mathop{cov}(\mu)\cdot\mathfrak{b}$.
\end{cor}
This may, or may not, place an additional value in $\{k(M) : \mathop{kc}(M)=\kappa\}$ - depending on whether $\lim_{\mu < \kappa} \mathop{cov}(\mu)\cdot\mathfrak{b}$ and $\mathop{cov}(\kappa)\cdot\mathfrak{b}$ are equal.


\paragraph{Powers - Attaining the Lower Bound} 

We next give a particular case when $k(M)$ realizes the lowest possible value, $kc(M)\cdot\mathfrak{b}$. 
\begin{thm}\label{th:powers} 

    Let $N$ be separable metrizable but not compact, and set $M=N^\omega$. Then $kc(M) \ge \mathfrak{d}$ and $k(M)=kc(M)$.     
    Further, if $|N|<\mathfrak{c}$ then $k(M)=kc(M)=\mathop{cov}(|N|)\cdot \mathfrak{d}$. 
\end{thm}
\begin{proof} 
    Since $N$ is not compact it contains a closed copy of $\omega$. Hence $M$ contains a closed copy of $\omega^\omega$, and so $kc(M) \ge kc(\omega^\omega)=\mathfrak{d}$. 
Again since $N$ is not compact, $M$ is not $\sigma$-compact, so $k(M) \ge kc(M)$. 
    We know from Theorem~3.6 of \cite{FG_Shape} that $\mathbf{kc}(M^\omega)\te\mathbf{kc}(M)^\omega \tq \mathbf{k}(M)$, so $kc(M^\omega) \ge k(M)$. 
    But $M^\omega$ is homeomorphic to $M$, from which it follows that $k(M)=kc(M)$. 


    Now suppose $|N|<\mathfrak{c}$, so compact subsets of $N$ are countable and $kc(N)=|N|$. We next show that $kc(M) \ge \mathop{cov}(|N|)$ by verifying $\textbf{kc}(M) \tq ([N]^\omega,\subseteq)$. 
    For $K$ in $\K(N^\omega)$ set $\phi_+(K)$ equal to $\bigcup \{ \pi_n(K) : n \in \omega\}$ if this latter set is infinite (and an arbitrary infinite subset of $N$ otherwise). 
    For $S=\{s_n : n \in \omega\}$ an infinite subset of $N$, set  $\psi_-(S)=(s_n)_n$. Then clearly, if $\psi_-(S) \in K$ then for every $n$, $s_n \in \pi_n(K)$, so $S \subseteq \phi_+(K)$. 

    To complete the argument we show $\mathop{cov}(|N|)\cdot\mathfrak{d} \ge kc(M)$. Let $\mathcal{C}$ be cofinal in $([N]^\omega,\subseteq)$ of size $\mathop{cov}(|N|)$. For each $C$ in $\mathcal{C}$ fix a bijective enumeration, $C=\{c_0,c_1,\ldots, c_n\ldots\}$. 
    Let $\mathcal{F}$ be cofinal in $(\omega^\omega,\le)$ of size $\mathfrak{d}$. 
    For each $C$ in $\mathcal{C}$ and $f$ from $\mathcal{F}$, define $K_{C,f}=\prod_n K_{C,f,n}$ where $K_{C,f,n}=\{c_0,c_1,\ldots,c_{f(n)}\}$. Note $K_{C,f}$ is a compact subset of $M=N^\omega$. 
    We check that the set $\{K_{C,f} : C \in \mathcal{C}, f \in \mathcal{F}\}$, which has size no more than $\mathop{cov}(|N|)\cdot\mathfrak{d}$, covers $M$. 

    Take any point $(x_n)_n$ in $N^\omega$. 
    Pick $C$ in $\mathcal{C}$ so that $\{x_n : n \in \omega\} \subseteq C=\{c_0,c_1,\ldots,c_i,\ldots\}$. For each $n$ let $g(n)$ be such that $x_n=c_{g(n)}$. Pick $f$ in $\mathcal{F}$ so that $g \le f$. Now, for each $n$, $x_n$ is in $K_{C,f,n}$. Hence $(x_n)_n$ is in $K_{C,f}$, as required. 
\end{proof}

A natural question is whether, in the above scenario where $M=N^\omega$, they have the same $k$-ness numbers: if $N$ is non-compact, separable metrizable then is $k(N^\omega)=k(N)$? 

\paragraph{The Number of Values Attained}

Recall that the $k$-spread at an uncountable $\kappa\le \mathfrak{c}$ is $\sigma_k(\kappa)=|\{k(M) : kc(M)=\kappa\}|$, namely, the number of values of $k$ attained for a fixed $kc$. We know from Theorem~\ref{th:k=sam}(i) that we always have $\sigma_k(\kappa) \ge 1$. 
Further, for every $n\in \mathbb{N}$ we know $\mathop{cov}(\aleph_n)=\aleph_n$, so $\sigma_k(\aleph_n)=1$. We now show that, modulo large cardinals, for any finite $K\ge 2$ there exist a model and a $\kappa$ with $\sigma_k(\kappa)=K$, and also a model and a $\kappa$ with $\sigma_k(\kappa)\ge\aleph_0$.

\begin{thm}\label{th:K_many}
Fix an integer $K\ge 2$. 
Let $\lambda=\aleph_{\omega\cdot K}$. 

Modulo large cardinals, it is consistent that $\mathfrak{b}=\aleph_1$, 
$\mathop{cov}(\lambda^+)=\lambda^{+K}$, and  
for $n=1,\ldots,K$ there are separable metrizable spaces, $M_n$, such that $kc(M_n)=\lambda^+$ while $k(M_n)=\lambda^{+n}$.

Hence in this model, $\kappa:=\lambda^+ \le \mathfrak{c}$ is uncountable and
\[\{k(M):kc(M)=\kappa\}=
[\lambda^+,\lambda^{+K}]=
[\kappa\cdot\mathfrak{b},\mathop{cov}(\kappa)\cdot\mathfrak{b}], \ \text{  so }\sigma_k(\kappa)=K.\] 
\end{thm}

\begin{thm}\label{th:ctbl_many}
Let $\lambda=\aleph_{\omega\cdot\omega}$. 

Modulo large cardinals, it is consistent that $\mathfrak{b}=\aleph_1$, 
$\mathop{cov}(\lambda^+)\ge\lambda^{+\omega}$, and  
for each integer $n\ge 1$ there are separable metrizable spaces, $M_n$, such that $kc(M_n)=\lambda^+$ while $k(M_n)=\lambda^{+n}$.

Hence in this model,  $\kappa:=\lambda^+ \le \mathfrak{c}$ is uncountable and
\[ [\lambda^+,\lambda^{+\omega})
\subseteq  \{k(M):kc(M)=\kappa\} \subseteq 
[\kappa\cdot\mathfrak{b},\mathop{cov}(\kappa)\cdot\mathfrak{b}], \ \text{  so }\sigma_k(\kappa)\ge\aleph_0.\] 
\end{thm}

In order to prove both these results we need to first construct, with the aid of large cardinals, models with suitable configurations of covering numbers, and then to find within those models the desired spaces, $M_n$. 
For the first step we rely on recent work of the second author and Cummings, \cite{GiltonCummings}.

\begin{thm}[Gilton \& Cummings, \cite{GiltonCummings}] Fix an integer $K \ge 2$. Modulo large cardinals, there is a model in which
(1) the \textsc{CH} holds, (2) for each $1\leq n\leq K$, $2^{\aleph_{\omega\cdot n}}=\aleph_{\omega\cdot K +n}$; and
    (3) for each $1\leq n\leq K$, $\aleph_{\omega\cdot n}$ admits a mod-finite scale of length $\aleph_{\omega\cdot K+n}$.
\end{thm}



\begin{proof}[Proof of  Theorem~\ref{th:K_many}]
Fix $K\ge 2$. 
Start with the model of Gilton and Cummings above. For each $1\leq n\leq K$, from (2), (3) together with Corollary~\ref{cor:ScalesIncreaseSam}   conclude that 
 $\mathop{cov}(\aleph_{\omega\cdot n})$ is equal to $\aleph_{\omega\cdot K+n}$. 
    Now do random real forcing to increase the value of $\mathfrak{c}$ strictly above $\aleph_{\omega\cdot K+K}$, while keeping $\mathfrak{b}$ and $\mathfrak{d}$ fixed at $\omega_1$. By Lemma~\ref{lem:cccPreservesCapturing}, this forcing doesn't change the relevant values of the covering numbers. 
Thus, in the resulting model, (a) $\mathfrak{b}=\mathfrak{d}=\omega_1$;
    (b) $\mathfrak{c}>\aleph_{\omega\cdot K+K}$; and
    (c) for each $1\leq n\leq K$, $\mathop{cov}(\aleph_{\omega\cdot n})=\aleph_{\omega\cdot K+n}$.

Let $\lambda=\aleph_{\omega\cdot K}$. 
Let $N$ be a subset of the reals of size $\aleph_\omega$, and set $M_\infty=N^\omega$. For each $1\le n \le K$, let $N_n$ have size $\aleph_{\omega\cdot n}$  but not be locally small, and set $M_n = M_\infty \oplus N_n$. 

We verify the claims of Theorem~\ref{th:K_many}. Item (a) simply asserts that $\mathfrak{b}=\omega_1$.
From Lemma~\ref{l:sam1}(iii) and from item (c), we see 
\[
\mathop{cov}(\lambda^+)=\mathop{cov}(\aleph_{\omega\cdot K+1})=\mathop{cov}(\aleph_{\omega\cdot K})\cdot \aleph_{\omega\cdot K+1}=\aleph_{\omega\cdot K+K}=\lambda^{+K}.
\]


The compact-covering number of $M_n$, $kc(M_n)$, is the larger of $kc(M_\infty)$ and $kc(N_n)$. 
From Theorem~\ref{th:powers} we know $kc(M_\infty)=\mathop{cov}(|N|)\cdot\mathfrak{d}=\mathop{cov}(\aleph_\omega)\cdot\aleph_1=\aleph_{\omega\cdot K+1}$. 
And since $|N_n|=\aleph_{\omega\cdot n} < \mathfrak{c}$ (by (b)), we have that $N_n$ is totally imperfect (since it has no compact sets of size $\mathfrak{c}$, on account of having size $<\mathfrak{c}$). Hence, 
$kc(N_n)=|N_n|=\aleph_{\omega\cdot n}$. 
Therefore $kc(M_n)=\aleph_{\omega\cdot K+1}=\lambda^+$, as required.

The $k$-ness number of $M_n$, $k(M_n)$, is the larger of $k(M_\infty)$ and $k(N_n)$. 
From Theorem~\ref{th:powers} we know 
$k(M_\infty)=kc(M_\infty)=\aleph_{\omega\cdot K+1}$. 
While from Theorem~\ref{th:k=sam}, as $kc(N_n)=\aleph_{\omega\cdot n}<\mathfrak{c}$ (by (b)), we have $k(N_n)=\mathop{cov}(\aleph_{\omega\cdot n})=\aleph_{\omega\cdot K+n}$. 
Hence $k(M_n)=\aleph_{\omega\cdot K+n}=\lambda^{+n}$, as required.    
\end{proof}

As input for the proof of   Theorem~\ref{th:ctbl_many} we utilize the following.
\begin{thm}[Gilton \& Cummings,\cite{GiltonCummings}] Modulo large cardinals, there is a model in which (1) the \textsc{CH} holds, 
    (2) for each $1\leq n$, $2^{\aleph_{\omega\cdot n}}=\aleph_{\omega^2 +n}$, and
    (3) for each $1\leq n$, 
    $\aleph_{\omega\cdot n}$ admits a mod-finite scale of length $\aleph_{\omega^2+n}$.
\end{thm}


\begin{proof}[Proof of  Theorem~\ref{th:ctbl_many}]
The argument is similar to that for Theorem~\ref{th:K_many} and we omit some details. Start with the model of Gilton and Cummings above. Applying a suitable random real forcing we get a model where (a) $\mathfrak{b}=\omega_1$, (b) $\mathfrak{c}>\aleph_{\omega^2+\omega}$, 
and
(c) $\mathop{cov}(\aleph_{\omega\cdot n})=\aleph_{\omega^2+n}$, for every $n \ge 1$.

Let $\lambda=\aleph_{\omega\cdot \omega}$. 
Let $N$ be a subset of the reals of size $\aleph_\omega$ and set $M_\infty=N^\omega$. For each $n \ge 1$ let $N_n$ have size $\aleph_{\omega\cdot n}$, but not be locally small, and set $M_n = M_\infty \oplus N_n$. 



Then, as above, in this model and with these spaces we have: $\mathfrak{b}=\aleph_1$, $\mathop{sam}(\lambda^+)\ge\aleph_{\omega\cdot\omega+\omega}=\lambda^{+\omega}$, and for every $n$, $kc(M_n)=\aleph_{\omega\cdot\omega+1}=\lambda^+$ while $k(M_n)=\aleph_{\omega\cdot\omega+n}=\lambda^{+n}$, all as required. 
\end{proof}


   
   
There is much that we do not know about the $k$-spread and the set $\{k(M) : \mathop{kc}(M)=\kappa\}$. 
We know, from the preceding result, it is consistent  that for $\kappa=\aleph_{\omega^2+1}$ the $k$-spread can be at least $\aleph_0$. 
Can the $k$-spread of some uncountable $\kappa \le \mathfrak{c}$ be uncountable? 
Can $\{k(M) : \mathop{kc}(M)=\kappa\}$  be a proper subset of $[\kappa\cdot\mathfrak{b},\mathop{cov}(\kappa)\cdot\mathfrak{b}]$? 






%XXX Is this combination of cardinals possible assuming sufficient large cardinals? TG: Not so sure. I misread the problem initially...

\paragraph{Van Douwen's Question}
In van Douwen's `Handbook' article \cite{vanDou} on small cardinals and topology he  asks the following: 
\begin{quote} Question 8.11: \ 
if $M$ is separable metrizable then is \\ \phantom{x}  $\cof(\K(M),\subseteq)=k(M)=\mathfrak{d}$ when $M$ is analytic, or at least Borel? 
\end{quote}
Evidently it is intended that $M$ not be $\sigma$-compact. 
Unasked by van Douwen, but also natural and considered in the literature, is the case when $M$ is co-analytic.


The answer is known for $\cof(\K(M),\subseteq)$. Under \textsc{CH} for any analytic $M$ we have $\cof(\K(M),\subseteq)=\mathfrak{d}=\aleph_1$. While Becker \cite{Becker} has constructed a model in which there is an analytic set $M$ with $\cof(\K(M),\subseteq)>\mathfrak{d}$. 
Van Engelen \cite{vanEng} showed that if $M$ is co-analytic but not locally compact then $\cof(\K(M),\subseteq) = \mathfrak{d}$. This also follows from earlier work of Fremlin, \cite{Fremlin}. 

It seems (see Vaughan \cite{Vau} and van Mill \cite{vanMill}) it was thought that these results completely answer  van Douwen's question. While that is true of $\cof(\K(M),\subseteq)$, it is not the case for $k(M)$. In fact our final theorem below shows that the situation with $k(M)$ is symmetrically opposite. \emph{In \textsc{ZFC}} the $k$-ness number of analytic, non $\sigma$-compact spaces is $\mathfrak{d}$. Hence for Becker's analytic $M$ we have $k(M)=\mathfrak{d}<\cof(\K(M),\subseteq)$. While it is consistent for there to be a co-analytic, non $\sigma$-compact space, $M$ say, whose $k$-ness number is strictly less than $\mathfrak{d}$. Combining this with the van Engelen/Fremlin result, for this co-analytic $M$ we have $k(M) <\mathfrak{d}=\cof(\K(M),\subseteq)$.  


\begin{thm}\label{th:VDsoln} \ 
\begin{enumerate}[noitemsep,topsep=1pt]
\item[(i)] If $M$ is analytic but not $\sigma$-compact, then  $k(M)=\mathfrak{d}$. 

\item[(ii)] It is consistent there is a co-analytic, non $\sigma$-compact $M$  with $\mathop{k}(M) < \mathfrak{d}$.
\end{enumerate}
\end{thm}
\begin{proof} 
For the first part, from Theorem~3.9 of \cite{FG_Shape}, we know that for analytic non $\sigma$-compact $M$ we have 
$\mathbf{k}(M) \te \omega^\omega$, and so, indeed, $k(M)=\cof(\omega^\omega)=\mathfrak{d}$. 


For the second claim we work in a model where $\aleph_1 < \mathfrak{p} \le \mathfrak{b} < \mathfrak{d}$,  and note that this is consistent.  Take any $\aleph_1$-sized subset, $M$ say, of the reals. 
Then $M$ is co-analytic because $|M|=\aleph_1 < \mathfrak{p}$. But $M$ is not $\sigma$-compact, in fact $\mathop{kc}(M)=\aleph_1$, because $|M|=\aleph_1<\mathfrak{c}$. Hence $\mathop{cov}(kc(M))\cdot \mathfrak{b}=\aleph_1\cdot\mathfrak{b}=\mathfrak{b}$ (Lemma~\ref{l:sam_aleph_n}), so  $\mathop{k}(M)=\mathfrak{b}$ (Theorem~\ref{th:k_sam}) and $\mathfrak{b} <\mathfrak{d}$ in our model. 
\end{proof}

\section{Acknowledgements} Thanks to Will Brian for his interest in this work, and for informing us \cite{Brian} about the equivalence of the sampling and covering numbers.

\begin{thebibliography}{99}

\bibitem{AbrahamMagidor}
U. Abraham, M. Magidor, \emph{Cardinal arithmetic}, in \emph{Handbook of Set Theory}, edited by Matthew Foreman and Akihiro Kanamori, Springer, Dordrecht, 2010, pp.~1149-1227.


\bibitem{Becker} H. Becker, \emph{Cofinal Families of Compact Subsets of an Analytic Set}, Proceedings of the American Mathematical Society
Vol. 106, No. 3 (Jul., 1989), pp. 853-856.

\bibitem{Brian} Personal Communication. Will Brian,  January 5th, 2026.


\bibitem{GiltonCummings} J. Cummings, T. Gilton, \emph{TBD}


\bibitem{vanDou} E. van Douwen, \emph{The integers and topology}, Handbook of set-theoretic
topology, pages 111-167, 1984.

\bibitem{vanEng} F. van Engelen, \emph{Cofinal families of compacta in separable metric spaces}, Proc. Amer.
Math. Soc. 104 (1988), 1271-1273.

\bibitem{FG_Shape_CC} Z. Feng,  P. Gartside,  \emph{The Shape of Compact Covers}. The Journal of Symbolic Logic (2024), pages 1-15. %10.1017/jsl.2024.53. 

\bibitem{FG_Shape} Z. Feng, P. Gartside, \emph{The Shape of Generating Families} (preprint)

\bibitem{Fremlin} D. Fremlin, \emph{Families of compact sets and Tukey’s ordering}, Atti Sem. Mat.
Fis. Univ. Modena 39, 1:29–50, 1991.

\bibitem{Gilton_PCF_Tukey} T. Gilton, 
\emph{\textsc{PCF} Theory and the Tukey Spectrum}, 
Fundamenta Mathematicae, 265 (2024), 15-33. 

\bibitem{Gitik:NegSCH} M. Gitik, \emph{The negation of the singular cardinal hypothesis from $o(\kappa)=\kappa^{++}$}, Ann. Pure and Appl. Logic, Vol 43, No. 3, 1989, pages 209-234.

\bibitem{Gitik:Strength} M. Gitik, \emph{The strength of the failure of the singular cardinal hypothesis}, Ann. Pure and Appl. Logic, Vol 51, No. 3, pages 215-240, 1991.


\bibitem{Gitik:noBounds}
M. Gitik. \emph{No bound for the first fixed point}. J. Math. Log, Vol 5, No. 2, pages 193-246, 2005.


\bibitem{Gitik:personal}
Personal Communication. Moti Gitik, by means of James Cummings, November 19th, 2025.

%\bibitem{JechCofinality} T. Jech, \emph{On the cofinality of countable products of cardinal numbers}, \emph{A Tribute to Paul Erd{\"o}s}, edited by A. Baker, B. Bollab{\'a}s, and A. Hajnal, Cambridge University Press, Cambridge, 1990, pp. 289-305.

\bibitem{Magidor}  M. Magidor, \emph{On the singular cardinals problem I}, Israel Journal of Mathematics, vol. 28 (1977), pp. 1-31.

\bibitem{vanMill} J. van Mill, \emph{Open Problems in van Douwen's Papers}, appeared in: Eric K. van Douwen,
\emph{Collected papers Vol. I, II.}, 
North-Holland Publishing Co., Amsterdam, 1994.


\bibitem{Sh1} S. Shelah, \emph{The singular cardinals problem; independence results}, Surveys in Set Theory (A. R. D. Mathias, editor), London Mathematical Society Lecture Note Series, vol. 87, Cambridge University
Press, Cambridge, 1983, pp. 116-134.

\bibitem{ShelahBook}
S. Shelah, \emph{Cardinal Arithmetic}, Oxford University Press, Oxford, 1994.

\bibitem{Vau} J. E. Vaughan, \emph{Small uncountable cardinals and topology}, Open Problems in Topology
(J. van Mill and G. M. Reed, eds.), North-Holland Publishing Company, Amsterdam,
1990, pp. 195-218.

\bibitem{Vojtas} P. Vojt\'{a}s, \emph{Generalized {G}alois-{T}ukey-connections between explicit relations on classical objects of real analysis}, Set theory of the reals, Israel Math. Conf. Proc., Bar-Ilan Univ., Ramat Gan, vol. 6, pages 619–643, 1993.


\end{thebibliography}


\end{document}


%To see this suppose $\mathcal{C}' \subseteq ([\kappa]^{<\omega})^{\omega}_{\infty}$ of size strictly less than $\kappa$. Then $S=\bigcup \{ \bigcup_n F_n : (F_n)_n \in \mathcal{C}'\}$ has size $<\kappa$, so pick $\alpha \in \kappa \setminus S$. Now we see $(\{\alpha\})_n$ witnesses that $\mathcal{C}'$ is not cofinal.


%From Lemma~\ref{l:sam1}(i) it follows $\mathop{seq}(M)=\lim_n \cof(P(\kappa_n))\cdot\kappa_n =\lim_{\mu <|M|} \mathop{sam}(\mu)\cdot \mathfrak{b}$. 



, where 
$<_\mathcal{U}$ is the total order on $\prod A_\infty$ induced by a free ultrafilter $\mathcal{U}$ on $A_\infty$. As $\mathcal{U}$ is free, clearly $\left(\prod A_\infty,<^*\right) \tq \left(\prod A_\infty,<_\mathcal{U}\right)$, and 
